% Unified short-paper replacement for princ.tex, mainRes.tex and exas.tex.
% Load only this file: % Unified short-paper replacement for princ.tex, mainRes.tex and exas.tex.
% Load only this file: % Unified short-paper replacement for princ.tex, mainRes.tex and exas.tex.
% Load only this file: % Unified short-paper replacement for princ.tex, mainRes.tex and exas.tex.
% Load only this file: \input{mainRes}.
% Keep the abstract's existing \ideaitem / \ideahead definitions unchanged.
% The six headings below are supplied by \ideahead, not copied manually.
\subsection{Main results}\label{s:mainRes}


The following enumeration follows the abstract; the text gives the main
formulas, their scope and relevant earlier results.
\begin{enumerate}
\ideahead{1}
The siphon identity and positivity give
\[
 J(0,y)=\begin{pmatrix}
 J_\Sigma^\perp(y)&0\\
 J_{\mathrm{mix},\Sigma}(y)&J_{\mathrm{tan},\Sigma}(y)
 \end{pmatrix},\qquad
 \begin{aligned}
 J_\Sigma^\perp&=\partial_zf_\Sigma(0,y),\\
 J_{\mathrm{mix},\Sigma}&=\partial_zf_{\Sigma^c}(0,y),\\
 J_{\mathrm{tan},\Sigma}&=\partial_yf_{\Sigma^c}(0,y).
 \end{aligned}
\]
The transverse block is Metzler even when the resident system is not
cooperative (Proposition~\ref{p:Tri}, \cite{AAHK,AA26}). For a regular
splitting $J_\Sigma^\perp=F_\Sigma-V_\Sigma$, with
$F_\Sigma\ge0$ and $V_\Sigma$ a nonsingular $M$-matrix,
\[
 s(J_\Sigma^\perp)<0
 \quad\Longleftrightarrow\quad
 R_\Sigma:=\rho(F_\Sigma V_\Sigma^{-1})<1
\]
is the next-generation matrix criterion \cite{Diek},\newline \cite{Van}.

For example, in the higher-order SIS model \eqref{eq:SIS}, $V_\Sigma=\operatorname{Diag}(\mu_i:i\in\Sigma)$.
For $i,j\in\Sigma$,
\[
 (F_\Sigma)_{ij}(y)=a_{ij}+
 \sum_{\substack{|E|\ge2\\E\cap\Sigma=\{j\}}}
 \beta_{iE}\prod_{k\in E\setminus\{j\}}y_k,
 \qquad J_\Sigma^\perp=F_\Sigma-V_\Sigma.
\]
Thus any term containing exactly one absent variable contributes to
invasion, regardless of its total degree. Terms vanishing together
with their first $z$-derivatives on the face leave both the resident
equations and the invasion matrix unchanged.
The scalar simplicial model \eqref{eq:simp} already exhibits bistability although its derivative at zero,
$\lambda-1$, is independent of $\lambda_\Delta$.

Theorem~\ref{t:siphMulti} constructs any prescribed finite number
of stable positive equilibria while preserving every siphon and every
transverse Jacobian $-V_\Sigma$. It uses nonnegative polynomials
$P_i$ and a smooth response vanishing on $(-\infty,0]$ and
positive on $(0,\infty)$. Its support condition is $V^{-1}c\gg0$,
where $c_i=b_i$ if $P_i\not\equiv0$ and $c_i=0$ otherwise.
The response depends on the prescribed count; no common slope bound
is imposed. The response is flat at onset, so all production derivatives on
siphon faces vanish: the transverse Jacobians are specifically
$-V_\Sigma$, rather than arbitrary prescribed matrices. This controls all boundary faces simultaneously, extending
the separation illustrated by the scalar contagion model. Siphon-based
persistence and boundary analysis have precedents in \cite{AdLS}.
The face property alone does not imply interior cooperativity:
$\dot z_1=-z_1z_2$, $\dot z_2=-z_2$ has a Metzler Jacobian at zero
but $\partial_{z_2}f_1=-z_1<0$ for $z_1>0$.

\ideahead{4}
For $\dot x=-Vx+\Phi(\beta Wx)$, assume a synaptic matrix $W\ge0$, $\beta\ge0$ and
that $\Phi\ge0$ is smooth on the attained arguments with $0\le \Phi'\le L$,
$L>0$. Theorem~\ref{t:KBsharp} compares
\[
 \bar R_0=\rho(\beta WV^{-1})
\]
with the reciprocal $1/L$ of the slope bound.
It extends the next-generation number: $J(x)=F(x)-V$ with
$F(x)=\beta\operatorname{Diag}(\Phi'(\beta Wx))W\le\beta LW$, so the classical number
$R(x_*)=\rho(F(x_*)V^{-1})$, which decides the stability of one equilibrium $x_*$ \cite{Van}, satisfies
$R(x_*)\le L\bar R_0$ for all equilibria and all responses in the class (Remark~\ref{r:NGMg}).
The transition is
\[
 \begin{cases}
 \bar R_0<\dfrac1L: & \parbox[t]{.7\linewidth}{every response in the class gives a unique globally exponentially attracting equilibrium, since $J(x)\le-V+\beta LW$;}\\[12pt]
 \bar R_0>\dfrac1L: & \parbox[t]{.7\linewidth}{for each $m\ge2$, some response in the class gives at least $m$ locally exponentially stable equilibria, including a choice with $\Phi>0$ and $\Phi'>0$ on $\mathbb R$ and all equilibria strictly positive.}
 \end{cases}
\]
No assertion at equality is needed. The response may depend on $m$;
this is not a count for one fixed response or polynomial degree.

Positive-feedback multistability is studied generally in
\cite{AngeliSontag2004}, \cite{AngeliFerrellSontag2004}.
The contraction below the bound is known: it is the case $B=I$
of the multivariable Lur'e theorem of \cite{Davydov}, for general $V$,
and \cite{Boud26} has the case $V=I$, $W\mathbf1=\mathbf1$; see also
\cite{Coogan2019,KawanoBesselinkCao2020}. The contribution here is its
constructive sharpness: with $V$, $W$, $L$ fixed and one common scalar response, every
prescribed finite stable count is attained above the bound, without irreducibility.

The construction principle, plateaus of the response on which $J=-V$, has a predecessor.
For $\dot u_i=-d_iu_i+\sum_jw_{ij}f_j(u_j)+I_i$ with piecewise linear
activations with $\ell+1$ plateaus, \cite[Theorem~3]{WangLuChen2010} give exactly
$(2\ell+1)^n$ equilibria, of which $(\ell+1)^n$ are locally stable, and estimate their
attraction basins; multilevel sigmoid activations are in their Section~5. Their condition (35)
requires self-feedback domination, which for a common activation implies
$w_{ii}>\ell\sum_{j\ne i}|w_{ij}|$. Thus it
cannot supply arbitrary counts on a fixed network with cross-coupling, whereas
Theorem~\ref{t:KBsharp} needs only $\bar R_0>1/L$, and asserts a lower bound rather than an exact count.
For example, take $V=I$, $W=\begin{pmatrix}0&1\\1&0\end{pmatrix}$ and $\beta L>1$. With
$y=\beta Wx$ the system becomes $\dot y=-y+\beta W\Phi(y)$, which has zero diagonal weights, so
(35) cannot hold, while Theorem~\ref{t:KBsharp} gives every finite count.
Theorem~\ref{t:siphMulti} uses the same mechanism; its addition is that $V$ and the
polynomials $P_i$ stay fixed, the siphon lattice and every transverse Jacobian are preserved,
and $V^{-1}c\gg0$ is the exact feasibility condition.

Betteti, Baggio, Bullo and Zampieri \cite{Betteti2025} construct symmetric firing-rate networks
whose prescribed two-level memory patterns are stable equilibria. Their formulas prescribe
response values at two argument levels, and their stability test controls the derivatives at those
levels, a precedent for the value-and-slope construction used here. Their synaptic matrix is
designed from the patterns, whereas Theorem~\ref{t:KBsharp} fixes the removal and synaptic matrices, dimension
and response slope bound and realizes every finite stable count by varying one common response.
Their construction does not impose the simultaneous preservation of siphons and transverse
matrices established in Theorem~\ref{t:siphMulti}. In their symmetric setting, an energy function
gives convergence to the set of equilibria. Its integral term is the convex generator of the
Fenchel--Bregman divergence of Subsection~\ref{s:FenBre}, which also recovers Volterra functions.

For an explicit three-population network, Theorem~\ref{t:KBsharp} is compatible with the
$2$-contraction criterion of \cite{OfirOvseevichMargaliot2024}. As a consequence of the two
results, arbitrarily many stable equilibria can coexist while every trajectory converges to an
equilibrium.

Equilibria of \eqref{eq:KB} are the fixed points of $T=V^{-1}\Pi$, and $DT(x_*)=V^{-1}\partial_x\Pi(x_*)$ is similar to the
next-generation matrix $\partial_x\Pi(x_*)V^{-1}$. For linear $P_i$, $T$ is a two-layer nonnegative neural network in the sense of
\cite{piotrowski2024fixed}. Their results apply when $T$ is monotone and weakly scalable, as for concave responses: the fixed point
is then unique and positive if $T$ is scalable, and the positive fixed points form an interval of a ray under a primitivity
condition. The responses of Theorems~\ref{t:KBsharp} and~\ref{t:siphMulti} and Proposition~\ref{p:Qua} are therefore not concave
where they are evaluated. Their certificate $T^m(0)\gg0$ excludes every nonempty siphon, since $T$ maps each siphon face into itself.

Linear-threshold networks have diagonal removal and clipped affine production, which is only locally Lipschitz.
\cite{Retnaraj2026} study a timescale family, with production $\tau^{-1}[Dx+\tau((W-D)x+h)]_0^1$, whose equilibrium set is independent
of $\tau$. Its invariant siphon faces are independent of $\tau$ as well: on $x_i=0$, the production vanishes exactly when
$((W-D)x+h)_i\le0$. Smooth bifurcation criteria require additional regularity at the clipping thresholds.

Unlike Theorem~\ref{t:siphMulti}, Theorem~\ref{t:KBsharp} fixes a slope bound
and linear $P_i$; its full-support construction has no nonempty
siphon. Neither theorem requires the other, and neither constructs
a cusp. The explicit cusp results use separate calculations and
continuation arguments.

\ifdefined\withSpatial
\ideahead{spat}
Add diffusion to \eqref{eq:KB}, $x_t=\mathcal D\Delta_\zeta x+\Pi(x)-Vx$ with $\zeta\in\mathbb R^n$, and linearize at a resident equilibrium on the face of a siphon
$\Sigma$. The absent variables obey $z_t=\mathcal D_\Sigma\Delta_\zeta z+J_\Sigma^\perp z$, with $J_\Sigma^\perp=F_\Sigma-V_\Sigma$, and the ansatz
$z=e^{-\lambda(e\cdot\zeta-ct)}q$, with $e$ a unit vector, gives the speed of exponential modes normal to $e$. In principle this speed depends on $e$; since the Laplacian assumes diffusion independent of the direction, it does not, and it is the speed of Proposition~\ref{t:spread}. The threshold is again next-generation: $c_*>0$ exactly when
$R_\Sigma>1$, and $c_*\sim2\sqrt{\langle\mathcal D\rangle s(J_\Sigma^\perp)}$ near it (Remark~\ref{r:movNGM}). The speed has the form of the
Freidlin--G\"artner formula. It concerns the linearized equation only: the flat-onset responses of Theorem~\ref{t:siphMulti} have
$R_\Sigma=0$ and no such speed, and with bistability a front can have positive speed although $R_\Sigma<1$.
For threshold responses, Theorem~\ref{t:SpdBd} carries Theorem~\ref{t:KBsharp} over to diffusion: below $1/L$ every solution converges to the
one equilibrium and there are no fronts, above it the $m$ stable states persist for every $\mathcal D$, and no response in $\mathcal C_L$ with $\Phi(0)=0$ spreads faster
than $\bar c(L)$, which we conjecture degree-one odds responses attain (Remark~\ref{r:SpdAtt}). For odds responses only $a_1$ enters the speed at a fixed resident
(Corollary~\ref{c:OddSpd}). The relay switches on two linear speeds, but numerically a front invades $E_{12}$ already while it is stable (Remark~\ref{r:RelSpd}), and the shared species $x_2$, absent from both residents, ignites where they meet (Remark~\ref{r:OvlIgn}, Figure~\ref{fig:overFront}). Every homogeneous Hurwitz equilibrium of a cooperative system stays stable under arbitrary diagonal diffusion, so all the constructed counts of stable equilibria survive, including the nine of Corollary~\ref{c:Nin} (Corollary~\ref{c:DifSta}); reductions with group-compatible diffusion, extinction certificates, global convergence on a box and, on bounded domains, homogeneous cusps extend as well (Corollary~\ref{c:DifAll}). Integrated Lyapunov functions acquire a diffusion term, which has a sign for Lyapunov--Volterra functions but not in general for network energies, and the composite basin certificate of Proposition~\ref{p:ComBas} extends after a change of storage (Proposition~\ref{p:DifLya}). Section~\ref{s:spaSta} recalls the Evans function and treats the stability of these fronts; above $c_*$ an exponential weight stabilizes the state ahead (Proposition~\ref{p:WeiSpe}).
\fi

\ideahead{3}
Partition the species into nonempty groups $C_1,\ldots,C_{n_G}$ and
let $H_{ia}=1$ for $i\in C_a$ and zero otherwise. Then $x=Hy$
assigns the common value $y_a$ to group $C_a$.
Assume $f(Hy;\theta)=H\widehat f(y;\theta)$ for all admissible
$y,\theta$. The group-constant states are invariant, and
\[
 J(Hy;\theta)H=H\widehat J(y;\theta),\qquad
 s(J(Hy;\theta))=s(\widehat J(y;\theta))
 \quad\text{if }J(Hy;\theta)\text{ is Metzler}.
\]
Theorem~\ref{t:Spe} transfers local exponential equilibrium stability
without irreducibility. A nondegenerate saddle-node or cusp also lifts
when the full Jacobian has a simple zero eigenvalue and all remaining
eigenvalues have negative real parts. These reductions support the
homogeneous and two-population calculations in
Theorems~\ref{t:KBhom} and~\ref{t:KBexchange}; the homogeneous cusp of a unimodal density
gives the hysteresis loop of Figure~\ref{fig:hyst}. For a symmetric response,
$\Phi(-q)=1-\Phi(q)$, Theorem~\ref{t:cuspAsym} gives the exact
asymmetry interval $|\eta|<\eta_*(\beta)$ on which both two-population cusps remain
admissible and nondegenerate.
Invariant synchrony subspaces and balanced partitions are established
in \cite{SGP03,golubitsky2005patterns}; here the cooperative spectral
equality supplies the stability test. Exact quotient reduction for row-stochastic
threshold networks, $WH=HB$, is in \cite{Boud26}, where a lifted equilibrium $Hy^*$ is
stable under the sufficient condition $\beta\max_aq_a<1$, with
$q_a=\Phi'(h+\beta(By^*)_a)$; the spectral equality replaces it by the exact test
$\beta\rho(\operatorname{Diag}(q)B)<1$. Invariance preserves initial
equality within groups, but does not imply synchronization from
arbitrary initial states.

\ideahead{odds}
For $\varphi(q)=\sum_{k=1}^m a_kq^k$, with $a_k\ge0$ and $a_m>0$,
put $\Phi(q)=\varphi(q)/(1+\varphi(q))$ for $q\ge0$ and $\Phi(q)=0$ for $q\le0$.
The zero set is fixed, and $\Phi'(0+)=a_1$.
For fixed $V$, $b_i>0$ and nonnegative polynomials $P_i$,
Proposition~\ref{p:Odd} preserves the siphon lattice and gives
\[
 J_\Sigma^\perp(y)=-V_\Sigma
 +a_1\operatorname{Diag}(b_i:i\in\Sigma)D_zP_\Sigma(0,y).
\]
Higher coefficients enter an invasion threshold only indirectly,
through the resident equilibrium. For $\dot x=-x+\phi(x)/(1+\phi(x))$,
with $\phi(0)=0$ and nonnegative polynomial coefficients, Proposition~\ref{p:Qua} combines
the exact minimum degree, four, for a positive nondegenerate cusp with a hyperbolically
stable origin, with an attaining construction.
Its reference polynomial is
\[
 \phi(q)=\frac{342q+368q^2+128q^3+1024q^4}{343}.
\]
For $\dot x=-x+\Phi(x)$, the equilibria $0$ and $1/2$ are hyperbolically
stable, while $1/8$ is a nondegenerate cusp with $f'''(1/8)=36/49$,
generically unfolded by $(a_1,a_2)$. More generally, the root parameters
$0<r<s$, $3r+s\le1$, give
\[
 c=[(1-r)^3(1-s)]^{-1},\qquad
 1-a_1=cr^3s,\qquad f'''(r)=6cr(s-r)(1-r).
\]
Within this family, the margin $1-a_1$ has supremum $1/81$, approached only as the cusp
degenerates. This bounds the uncoupled origin's stability margin,
not the coupling range of a continued cusp.
Theorem~\ref{t:Net} continues the cusp and the two stable equilibria
under coupling that preserves singleton siphons and a transversely
stable boundary equilibrium.

Bistability and hysteresis under rational genetic feedback are established:
\cite{Hasty2001} derive
\begin{equation}\label{eq:Hasty}
 \dot x=m\frac{1+x^2+\alpha\sigma_1x^4}{1+x^2+\sigma_1x^4+\sigma_1\sigma_2x^6}-\gamma_xx
\end{equation}
by eliminating fast dimerization and promoter-binding reactions, with hysteresis as the degradation
rate varies; see also \cite{Isaacs2003}. Its production at zero is $m>0$, so the origin is not an equilibrium, and its numerator
and denominator are different polynomials; its $x^4$ term therefore has a different role from
the degree of $\phi$ here.
Earlier biological feedback models already exhibit cusps associated
with rational response functions. In particular, van Hoek and
Hogeweg \cite[Eq.~(21)]{vanHoek} locate a cusp by imposing
a triple root on the polynomial equilibrium equation.
Izhikevich \cite[Theorem~1]{Izhikevich1998} derives coupled cubic
normal forms for weakly coupled systems whose uncoupled components
already have cusps. These results precede our use of polynomial
equilibrium equations and coupling near a cusp.

We have not found the minimum-degree result in the preceding
literature; it concerns the displayed odds family, not rational genetic responses in general. The local cusp criterion is standard \cite{kuznetsov2004};
the additional result is its explicit realization subject to these
positivity and boundary constraints.

\ideahead{overl}
Responses positive at every finite argument, such as Gaussian and logistic distribution functions,
exclude nonempty siphons (Remark~\ref{r:KBnosiph}). For zero-onset responses,
Proposition~\ref{p:InvF} gives the criterion:
for each siphon $\Sigma$, every production monomial for a species
$i\in\Sigma$ must contain at least one variable from $\Sigma$,
and $V_{\Sigma\Sigma^c}=0$. For a singleton, this forces its own
variable to divide each production monomial (Remark~\ref{r:Mul}).
For linear $P_i$, minimal siphons are the source strongly connected components of the
graph of $W$ and $V$ (Remark~\ref{r:KBsiphgraph}); for higher-order $P_i$ they can overlap
(Example~\ref{e:KBoverlap}). The degree of $P_i$ in the absent variables and the onset order
of $\Phi_i$ fix the order of production along transverse rays (Remark~\ref{r:KBtransOrder}).
The three-population network
\[
 \begin{aligned}
 \dot x_1&=-x_1+\Phi_a(x_1+\alpha x_2),\\
 \dot x_2&=-x_2+\Phi_a(x_2+\beta x_1x_3),\\
 \dot x_3&=-x_3+\Phi_a(x_3+\gamma x_2)
 \end{aligned}
\]
with positive couplings has overlapping minimal siphons
$\Sigma_{12}=\{1,2\}$ and $\Sigma_{23}=\{2,3\}$.
Their faces carry $(0,0,s_a)$ and $(s_a,0,0)$, where $s_a>0$
is the continued stable scalar equilibrium. Their invasion numbers are
\[
 R_{\Sigma_{12}}=a_1(1+\sqrt{\alpha\beta s_a}),\qquad
 R_{\Sigma_{23}}=a_1(1+\sqrt{\gamma\beta s_a}).
\]
They differ when $\alpha\ne\gamma$ and depend on the higher odds
coefficients through $s_a$.
These siphons share species $2$, whose production term $\beta x_1x_3$ enters both transverse matrices, while
both residents are independent of $\beta$. This \textbf{double boundary transcritical relay} takes the two residents
through transcritical bifurcations in turn by varying $\beta$ alone (Remark~\ref{r:KBrelay}, Figure~\ref{fig:relay}), which cannot
happen with disjoint siphons.
For sufficiently small positive couplings, Theorem~\ref{t:KBoverlapCusp}
constructs eight hyperbolic stable equilibria coexisting with an
interior cusp. Its positive cubic coefficient makes the middle
branch of the local unfolding stable, yielding at least nine stable
equilibria in a nearby open parameter region (Corollary~\ref{c:Nin}).
The cusp equilibrium itself is not counted among the eight.
A composite function, the scalar quartic energy on the resident coordinate plus a Perron-weighted
linear term on the absent coordinates, certifies an explicit subset of the basin of attraction of each boundary resident
(Proposition~\ref{p:ComBas}).
Earlier weak-coupling constructions \cite{Izhikevich1998,Boud26}
provide precedents for multiple stable states; the additional
constraints here preserve overlapping siphons and distinct explicit
boundary invasion thresholds.

\ideahead{5}
If the absent variables satisfy
\[
 \dot x_\Sigma\le(C_\Sigma-V_\Sigma)x_\Sigma,\qquad
 C_\Sigma\ge0,\qquad \rho(C_\Sigma V_\Sigma^{-1})<1,
\]
with $V_\Sigma$ a nonsingular $M$-matrix, then a constant weight
$w\gg0$ with $(C_\Sigma-V_\Sigma)^\top w\ll0$ certifies exponential
extinction (Theorem~\ref{t:Cmp}).
For a cooperative $C^2$ system on a compact invariant box with a
Hurwitz equilibrium $e$, Theorem~\ref{t:EqWgt} characterizes global
convergence to $e$ by a $C^1$ weight $v\gg0$ and $c>0$ satisfying
\[
 (\partial_xv)f+J^\top v\le-cv.
\]
The construction uses the variational flow
$\Psi(t,x)=\partial_x\varphi_t(x)$ of $f$, where
$\partial_t\Psi=J(\varphi_t(x))\Psi$ and $\Psi(0,x)=I$:
\[
 v_T(x)=\int_0^T\Psi(t,x)^\top q\,dt,\qquad q\gg0,
\]
for a sufficiently large finite $T$.
The contracting distance is
\[
 d_v(x,y)=\inf_{\gamma(0)=x,\ \gamma(1)=y}
 \int_0^1 v(\gamma(\tau))^\top|\gamma'(\tau)|\,d\tau,
\]
where the infimum is over piecewise $C^1$ paths in the box.
It need not equal $v(x)^\top|x-y|$.
Theorem~\ref{t:Wgt} gives the corresponding positive-weight converse
for Metzler variational equations over a compact flow.

The sufficient direction is established. For Metzler $J$, the inequality above is the
weighted $\ell^1$ case, $\Theta=\operatorname{Diag}(v)$, of Coogan's contraction condition
$\mu(\dot\Theta\Theta^{-1}+\Theta J\Theta^{-1})\le-c$ \cite{Coogan2019}; his separable Lyapunov
theorems and sum-of-squares search use weights $\theta_i(x_i)$.
The converse is also established for monotone systems: on a convex, forward and backward
invariant set, with bounded Jacobian, Kawano, Besselink and Cao
\cite[Theorem~4.3, Corollary~4.8]{KawanoBesselinkCao2020} prove that incremental exponential
stability is equivalent to a bounded, uniformly positive $C^1$ weight $v(x)$ satisfying this
inequality, and their proof uses the finite-horizon weight above
\cite[Appendix~B, (61)]{KawanoBesselinkCao2020}.
Theorem~\ref{t:Wgt} formulates this principle over a compact base flow: the base need not be
cooperative, $A$ need not be its Jacobian, and only forward evolution is used.
Theorem~\ref{t:EqWgt} is a compact-domain consequence, starting from global convergence to a
Hurwitz equilibrium on a forward invariant set. Existence of a polynomial weight on a
rectangular box is an approximation corollary. Separability here is in the tangent variables,
with $v_i$ depending on the full state.
Related converse results are in
\cite{LohmillerSlotine1998,ForniSepulchre2014,Giesl2015}.
The weights here also cover extinction. These results describe global convergence;
they are not prerequisites for the cusp constructions. Multistability
excludes contraction on a box containing the distinct equilibria.

Lyapunov--Volterra functions fit the same picture.
The function $\sum_ic_i(x_i-x_i^*-x_i^*\log(x_i/x_i^*))$ is a weighted sum of Fenchel--Bregman
divergences of the convex generator $x\log x-x$ of the exponential response
(Subsection~\ref{s:FenBre}). Its zero-reference limit gives the linear terms of the extinction
certificate of Theorem~\ref{t:Cmp}. The inverse-response primitive also enters the
response-integral energy of symmetric threshold networks (Proposition~\ref{p:SymEne}) and of the
scalar quartic odds equation. These divergences are nonnegative; the network
energy additionally contains a quadratic coupling term. Its decrease, and its local minimum at a
stable equilibrium, follow from the separate calculations given below.
Proposition~\ref{p:ComLya} combines a resident storage with such a linear term, and Proposition~\ref{p:LVboundary} applies this with a Lyapunov--Volterra storage to the three-population network.
\end{enumerate}

% \paragraph{Odds-polynomial responses and the minimum cusp degree.}
% Autoregulatory genetic models already combine saturating production,
% positive feedback and bistability; Isaacs et al.~\cite{Isaacs2003}
% provide a theoretical and experimental example.
% Multiple cusp bifurcations and multiple attracting states in weakly
% connected neural networks were studied by Izhikevich
% \cite{Izhikevich1998}.
% The specific contribution of
% Proposition~\ref{p:Qua}  is the exact minimum degree within the stated
% odds-polynomial class: degree four is necessary and sufficient for a
% stable origin and a positive nondegenerate cusp to coexist.
% The response zero set fixes the siphons, its onset slope determines
% the transverse production derivatives at a fixed resident state,
% and higher coefficients control the interior equilibrium geometry.
% Theorem~\ref{t:Net} then continues the cusp under coupling that
% preserves the prescribed boundary faces.
%
% \paragraph{Overlapping siphons and coexistence of stable states.}
% Persistence of hyperbolic equilibria under small perturbations is
% standard, and weakly coupled multistable networks already occur in
% \cite{Izhikevich1998}. Theorem~\ref{t:KBoverlapCusp} imposes additional
% boundary constraints on such a construction.
% The minimal siphons $\{1,2\}$ and $\{2,3\}$ remain fixed, their proper
% faces carry nonzero resident equilibria with distinct invasion numbers,
% and eight hyperbolic stable equilibria coexist with an interior cusp.
% Corollary~\ref{c:Nin} obtains a ninth stable equilibrium from its
% unfolding. This is an explicit realization of the general
% siphon-preservation rule, rather than a claim that weak coupling
% or multistability itself is new.
%
% \paragraph{Positive weights and converse contraction.}
% Coogan \cite{Coogan2019} develops contraction criteria using
% state-dependent weighted norms, while Kawano, Besselink and Cao
% \cite{KawanoBesselinkCao2020} prove converse results for monotone
% systems and positive linear time-varying systems.
% Theorem~\ref{t:EqWgt} belongs to this established line of work.
% For a cooperative system on a compact invariant box, it constructs
% positive contraction weights from global convergence to a Hurwitz
% equilibrium. Its role here is to turn convergence, including
% extinction, into an explicit variational-flow certificate.
% The siphon framework identifies the relevant boundary subsystems;
% the converse contraction principle has an independent literature..
% Keep the abstract's existing \ideaitem / \ideahead definitions unchanged.
% The six headings below are supplied by \ideahead, not copied manually.
\subsection{Main results}\label{s:mainRes}


The following enumeration follows the abstract; the text gives the main
formulas, their scope and relevant earlier results.
\begin{enumerate}
\ideahead{1}
The siphon identity and positivity give
\[
 J(0,y)=\begin{pmatrix}
 J_\Sigma^\perp(y)&0\\
 J_{\mathrm{mix},\Sigma}(y)&J_{\mathrm{tan},\Sigma}(y)
 \end{pmatrix},\qquad
 \begin{aligned}
 J_\Sigma^\perp&=\partial_zf_\Sigma(0,y),\\
 J_{\mathrm{mix},\Sigma}&=\partial_zf_{\Sigma^c}(0,y),\\
 J_{\mathrm{tan},\Sigma}&=\partial_yf_{\Sigma^c}(0,y).
 \end{aligned}
\]
The transverse block is Metzler even when the resident system is not
cooperative (Proposition~\ref{p:Tri}, \cite{AAHK,AA26}). For a regular
splitting $J_\Sigma^\perp=F_\Sigma-V_\Sigma$, with
$F_\Sigma\ge0$ and $V_\Sigma$ a nonsingular $M$-matrix,
\[
 s(J_\Sigma^\perp)<0
 \quad\Longleftrightarrow\quad
 R_\Sigma:=\rho(F_\Sigma V_\Sigma^{-1})<1
\]
is the next-generation matrix criterion \cite{Diek},\newline \cite{Van}.

For example, in the higher-order SIS model \eqref{eq:SIS}, $V_\Sigma=\operatorname{Diag}(\mu_i:i\in\Sigma)$.
For $i,j\in\Sigma$,
\[
 (F_\Sigma)_{ij}(y)=a_{ij}+
 \sum_{\substack{|E|\ge2\\E\cap\Sigma=\{j\}}}
 \beta_{iE}\prod_{k\in E\setminus\{j\}}y_k,
 \qquad J_\Sigma^\perp=F_\Sigma-V_\Sigma.
\]
Thus any term containing exactly one absent variable contributes to
invasion, regardless of its total degree. Terms vanishing together
with their first $z$-derivatives on the face leave both the resident
equations and the invasion matrix unchanged.
The scalar simplicial model \eqref{eq:simp} already exhibits bistability although its derivative at zero,
$\lambda-1$, is independent of $\lambda_\Delta$.

Theorem~\ref{t:siphMulti} constructs any prescribed finite number
of stable positive equilibria while preserving every siphon and every
transverse Jacobian $-V_\Sigma$. It uses nonnegative polynomials
$P_i$ and a smooth response vanishing on $(-\infty,0]$ and
positive on $(0,\infty)$. Its support condition is $V^{-1}c\gg0$,
where $c_i=b_i$ if $P_i\not\equiv0$ and $c_i=0$ otherwise.
The response depends on the prescribed count; no common slope bound
is imposed. The response is flat at onset, so all production derivatives on
siphon faces vanish: the transverse Jacobians are specifically
$-V_\Sigma$, rather than arbitrary prescribed matrices. This controls all boundary faces simultaneously, extending
the separation illustrated by the scalar contagion model. Siphon-based
persistence and boundary analysis have precedents in \cite{AdLS}.
The face property alone does not imply interior cooperativity:
$\dot z_1=-z_1z_2$, $\dot z_2=-z_2$ has a Metzler Jacobian at zero
but $\partial_{z_2}f_1=-z_1<0$ for $z_1>0$.

\ideahead{4}
For $\dot x=-Vx+\Phi(\beta Wx)$, assume a synaptic matrix $W\ge0$, $\beta\ge0$ and
that $\Phi\ge0$ is smooth on the attained arguments with $0\le \Phi'\le L$,
$L>0$. Theorem~\ref{t:KBsharp} compares
\[
 \bar R_0=\rho(\beta WV^{-1})
\]
with the reciprocal $1/L$ of the slope bound.
It extends the next-generation number: $J(x)=F(x)-V$ with
$F(x)=\beta\operatorname{Diag}(\Phi'(\beta Wx))W\le\beta LW$, so the classical number
$R(x_*)=\rho(F(x_*)V^{-1})$, which decides the stability of one equilibrium $x_*$ \cite{Van}, satisfies
$R(x_*)\le L\bar R_0$ for all equilibria and all responses in the class (Remark~\ref{r:NGMg}).
The transition is
\[
 \begin{cases}
 \bar R_0<\dfrac1L: & \parbox[t]{.7\linewidth}{every response in the class gives a unique globally exponentially attracting equilibrium, since $J(x)\le-V+\beta LW$;}\\[12pt]
 \bar R_0>\dfrac1L: & \parbox[t]{.7\linewidth}{for each $m\ge2$, some response in the class gives at least $m$ locally exponentially stable equilibria, including a choice with $\Phi>0$ and $\Phi'>0$ on $\mathbb R$ and all equilibria strictly positive.}
 \end{cases}
\]
No assertion at equality is needed. The response may depend on $m$;
this is not a count for one fixed response or polynomial degree.

Positive-feedback multistability is studied generally in
\cite{AngeliSontag2004}, \cite{AngeliFerrellSontag2004}.
The contraction below the bound is known: it is the case $B=I$
of the multivariable Lur'e theorem of \cite{Davydov}, for general $V$,
and \cite{Boud26} has the case $V=I$, $W\mathbf1=\mathbf1$; see also
\cite{Coogan2019,KawanoBesselinkCao2020}. The contribution here is its
constructive sharpness: with $V$, $W$, $L$ fixed and one common scalar response, every
prescribed finite stable count is attained above the bound, without irreducibility.

The construction principle, plateaus of the response on which $J=-V$, has a predecessor.
For $\dot u_i=-d_iu_i+\sum_jw_{ij}f_j(u_j)+I_i$ with piecewise linear
activations with $\ell+1$ plateaus, \cite[Theorem~3]{WangLuChen2010} give exactly
$(2\ell+1)^n$ equilibria, of which $(\ell+1)^n$ are locally stable, and estimate their
attraction basins; multilevel sigmoid activations are in their Section~5. Their condition (35)
requires self-feedback domination, which for a common activation implies
$w_{ii}>\ell\sum_{j\ne i}|w_{ij}|$. Thus it
cannot supply arbitrary counts on a fixed network with cross-coupling, whereas
Theorem~\ref{t:KBsharp} needs only $\bar R_0>1/L$, and asserts a lower bound rather than an exact count.
For example, take $V=I$, $W=\begin{pmatrix}0&1\\1&0\end{pmatrix}$ and $\beta L>1$. With
$y=\beta Wx$ the system becomes $\dot y=-y+\beta W\Phi(y)$, which has zero diagonal weights, so
(35) cannot hold, while Theorem~\ref{t:KBsharp} gives every finite count.
Theorem~\ref{t:siphMulti} uses the same mechanism; its addition is that $V$ and the
polynomials $P_i$ stay fixed, the siphon lattice and every transverse Jacobian are preserved,
and $V^{-1}c\gg0$ is the exact feasibility condition.

Betteti, Baggio, Bullo and Zampieri \cite{Betteti2025} construct symmetric firing-rate networks
whose prescribed two-level memory patterns are stable equilibria. Their formulas prescribe
response values at two argument levels, and their stability test controls the derivatives at those
levels, a precedent for the value-and-slope construction used here. Their synaptic matrix is
designed from the patterns, whereas Theorem~\ref{t:KBsharp} fixes the removal and synaptic matrices, dimension
and response slope bound and realizes every finite stable count by varying one common response.
Their construction does not impose the simultaneous preservation of siphons and transverse
matrices established in Theorem~\ref{t:siphMulti}. In their symmetric setting, an energy function
gives convergence to the set of equilibria. Its integral term is the convex generator of the
Fenchel--Bregman divergence of Subsection~\ref{s:FenBre}, which also recovers Volterra functions.

For an explicit three-population network, Theorem~\ref{t:KBsharp} is compatible with the
$2$-contraction criterion of \cite{OfirOvseevichMargaliot2024}. As a consequence of the two
results, arbitrarily many stable equilibria can coexist while every trajectory converges to an
equilibrium.

Equilibria of \eqref{eq:KB} are the fixed points of $T=V^{-1}\Pi$, and $DT(x_*)=V^{-1}\partial_x\Pi(x_*)$ is similar to the
next-generation matrix $\partial_x\Pi(x_*)V^{-1}$. For linear $P_i$, $T$ is a two-layer nonnegative neural network in the sense of
\cite{piotrowski2024fixed}. Their results apply when $T$ is monotone and weakly scalable, as for concave responses: the fixed point
is then unique and positive if $T$ is scalable, and the positive fixed points form an interval of a ray under a primitivity
condition. The responses of Theorems~\ref{t:KBsharp} and~\ref{t:siphMulti} and Proposition~\ref{p:Qua} are therefore not concave
where they are evaluated. Their certificate $T^m(0)\gg0$ excludes every nonempty siphon, since $T$ maps each siphon face into itself.

Linear-threshold networks have diagonal removal and clipped affine production, which is only locally Lipschitz.
\cite{Retnaraj2026} study a timescale family, with production $\tau^{-1}[Dx+\tau((W-D)x+h)]_0^1$, whose equilibrium set is independent
of $\tau$. Its invariant siphon faces are independent of $\tau$ as well: on $x_i=0$, the production vanishes exactly when
$((W-D)x+h)_i\le0$. Smooth bifurcation criteria require additional regularity at the clipping thresholds.

Unlike Theorem~\ref{t:siphMulti}, Theorem~\ref{t:KBsharp} fixes a slope bound
and linear $P_i$; its full-support construction has no nonempty
siphon. Neither theorem requires the other, and neither constructs
a cusp. The explicit cusp results use separate calculations and
continuation arguments.

\ifdefined\withSpatial
\ideahead{spat}
Add diffusion to \eqref{eq:KB}, $x_t=\mathcal D\Delta_\zeta x+\Pi(x)-Vx$ with $\zeta\in\mathbb R^n$, and linearize at a resident equilibrium on the face of a siphon
$\Sigma$. The absent variables obey $z_t=\mathcal D_\Sigma\Delta_\zeta z+J_\Sigma^\perp z$, with $J_\Sigma^\perp=F_\Sigma-V_\Sigma$, and the ansatz
$z=e^{-\lambda(e\cdot\zeta-ct)}q$, with $e$ a unit vector, gives the speed of exponential modes normal to $e$. In principle this speed depends on $e$; since the Laplacian assumes diffusion independent of the direction, it does not, and it is the speed of Proposition~\ref{t:spread}. The threshold is again next-generation: $c_*>0$ exactly when
$R_\Sigma>1$, and $c_*\sim2\sqrt{\langle\mathcal D\rangle s(J_\Sigma^\perp)}$ near it (Remark~\ref{r:movNGM}). The speed has the form of the
Freidlin--G\"artner formula. It concerns the linearized equation only: the flat-onset responses of Theorem~\ref{t:siphMulti} have
$R_\Sigma=0$ and no such speed, and with bistability a front can have positive speed although $R_\Sigma<1$.
For threshold responses, Theorem~\ref{t:SpdBd} carries Theorem~\ref{t:KBsharp} over to diffusion: below $1/L$ every solution converges to the
one equilibrium and there are no fronts, above it the $m$ stable states persist for every $\mathcal D$, and no response in $\mathcal C_L$ with $\Phi(0)=0$ spreads faster
than $\bar c(L)$, which we conjecture degree-one odds responses attain (Remark~\ref{r:SpdAtt}). For odds responses only $a_1$ enters the speed at a fixed resident
(Corollary~\ref{c:OddSpd}). The relay switches on two linear speeds, but numerically a front invades $E_{12}$ already while it is stable (Remark~\ref{r:RelSpd}), and the shared species $x_2$, absent from both residents, ignites where they meet (Remark~\ref{r:OvlIgn}, Figure~\ref{fig:overFront}). Every homogeneous Hurwitz equilibrium of a cooperative system stays stable under arbitrary diagonal diffusion, so all the constructed counts of stable equilibria survive, including the nine of Corollary~\ref{c:Nin} (Corollary~\ref{c:DifSta}); reductions with group-compatible diffusion, extinction certificates, global convergence on a box and, on bounded domains, homogeneous cusps extend as well (Corollary~\ref{c:DifAll}). Integrated Lyapunov functions acquire a diffusion term, which has a sign for Lyapunov--Volterra functions but not in general for network energies, and the composite basin certificate of Proposition~\ref{p:ComBas} extends after a change of storage (Proposition~\ref{p:DifLya}). Section~\ref{s:spaSta} recalls the Evans function and treats the stability of these fronts; above $c_*$ an exponential weight stabilizes the state ahead (Proposition~\ref{p:WeiSpe}).
\fi

\ideahead{3}
Partition the species into nonempty groups $C_1,\ldots,C_{n_G}$ and
let $H_{ia}=1$ for $i\in C_a$ and zero otherwise. Then $x=Hy$
assigns the common value $y_a$ to group $C_a$.
Assume $f(Hy;\theta)=H\widehat f(y;\theta)$ for all admissible
$y,\theta$. The group-constant states are invariant, and
\[
 J(Hy;\theta)H=H\widehat J(y;\theta),\qquad
 s(J(Hy;\theta))=s(\widehat J(y;\theta))
 \quad\text{if }J(Hy;\theta)\text{ is Metzler}.
\]
Theorem~\ref{t:Spe} transfers local exponential equilibrium stability
without irreducibility. A nondegenerate saddle-node or cusp also lifts
when the full Jacobian has a simple zero eigenvalue and all remaining
eigenvalues have negative real parts. These reductions support the
homogeneous and two-population calculations in
Theorems~\ref{t:KBhom} and~\ref{t:KBexchange}; the homogeneous cusp of a unimodal density
gives the hysteresis loop of Figure~\ref{fig:hyst}. For a symmetric response,
$\Phi(-q)=1-\Phi(q)$, Theorem~\ref{t:cuspAsym} gives the exact
asymmetry interval $|\eta|<\eta_*(\beta)$ on which both two-population cusps remain
admissible and nondegenerate.
Invariant synchrony subspaces and balanced partitions are established
in \cite{SGP03,golubitsky2005patterns}; here the cooperative spectral
equality supplies the stability test. Exact quotient reduction for row-stochastic
threshold networks, $WH=HB$, is in \cite{Boud26}, where a lifted equilibrium $Hy^*$ is
stable under the sufficient condition $\beta\max_aq_a<1$, with
$q_a=\Phi'(h+\beta(By^*)_a)$; the spectral equality replaces it by the exact test
$\beta\rho(\operatorname{Diag}(q)B)<1$. Invariance preserves initial
equality within groups, but does not imply synchronization from
arbitrary initial states.

\ideahead{odds}
For $\varphi(q)=\sum_{k=1}^m a_kq^k$, with $a_k\ge0$ and $a_m>0$,
put $\Phi(q)=\varphi(q)/(1+\varphi(q))$ for $q\ge0$ and $\Phi(q)=0$ for $q\le0$.
The zero set is fixed, and $\Phi'(0+)=a_1$.
For fixed $V$, $b_i>0$ and nonnegative polynomials $P_i$,
Proposition~\ref{p:Odd} preserves the siphon lattice and gives
\[
 J_\Sigma^\perp(y)=-V_\Sigma
 +a_1\operatorname{Diag}(b_i:i\in\Sigma)D_zP_\Sigma(0,y).
\]
Higher coefficients enter an invasion threshold only indirectly,
through the resident equilibrium. For $\dot x=-x+\phi(x)/(1+\phi(x))$,
with $\phi(0)=0$ and nonnegative polynomial coefficients, Proposition~\ref{p:Qua} combines
the exact minimum degree, four, for a positive nondegenerate cusp with a hyperbolically
stable origin, with an attaining construction.
Its reference polynomial is
\[
 \phi(q)=\frac{342q+368q^2+128q^3+1024q^4}{343}.
\]
For $\dot x=-x+\Phi(x)$, the equilibria $0$ and $1/2$ are hyperbolically
stable, while $1/8$ is a nondegenerate cusp with $f'''(1/8)=36/49$,
generically unfolded by $(a_1,a_2)$. More generally, the root parameters
$0<r<s$, $3r+s\le1$, give
\[
 c=[(1-r)^3(1-s)]^{-1},\qquad
 1-a_1=cr^3s,\qquad f'''(r)=6cr(s-r)(1-r).
\]
Within this family, the margin $1-a_1$ has supremum $1/81$, approached only as the cusp
degenerates. This bounds the uncoupled origin's stability margin,
not the coupling range of a continued cusp.
Theorem~\ref{t:Net} continues the cusp and the two stable equilibria
under coupling that preserves singleton siphons and a transversely
stable boundary equilibrium.

Bistability and hysteresis under rational genetic feedback are established:
\cite{Hasty2001} derive
\begin{equation}\label{eq:Hasty}
 \dot x=m\frac{1+x^2+\alpha\sigma_1x^4}{1+x^2+\sigma_1x^4+\sigma_1\sigma_2x^6}-\gamma_xx
\end{equation}
by eliminating fast dimerization and promoter-binding reactions, with hysteresis as the degradation
rate varies; see also \cite{Isaacs2003}. Its production at zero is $m>0$, so the origin is not an equilibrium, and its numerator
and denominator are different polynomials; its $x^4$ term therefore has a different role from
the degree of $\phi$ here.
Earlier biological feedback models already exhibit cusps associated
with rational response functions. In particular, van Hoek and
Hogeweg \cite[Eq.~(21)]{vanHoek} locate a cusp by imposing
a triple root on the polynomial equilibrium equation.
Izhikevich \cite[Theorem~1]{Izhikevich1998} derives coupled cubic
normal forms for weakly coupled systems whose uncoupled components
already have cusps. These results precede our use of polynomial
equilibrium equations and coupling near a cusp.

We have not found the minimum-degree result in the preceding
literature; it concerns the displayed odds family, not rational genetic responses in general. The local cusp criterion is standard \cite{kuznetsov2004};
the additional result is its explicit realization subject to these
positivity and boundary constraints.

\ideahead{overl}
Responses positive at every finite argument, such as Gaussian and logistic distribution functions,
exclude nonempty siphons (Remark~\ref{r:KBnosiph}). For zero-onset responses,
Proposition~\ref{p:InvF} gives the criterion:
for each siphon $\Sigma$, every production monomial for a species
$i\in\Sigma$ must contain at least one variable from $\Sigma$,
and $V_{\Sigma\Sigma^c}=0$. For a singleton, this forces its own
variable to divide each production monomial (Remark~\ref{r:Mul}).
For linear $P_i$, minimal siphons are the source strongly connected components of the
graph of $W$ and $V$ (Remark~\ref{r:KBsiphgraph}); for higher-order $P_i$ they can overlap
(Example~\ref{e:KBoverlap}). The degree of $P_i$ in the absent variables and the onset order
of $\Phi_i$ fix the order of production along transverse rays (Remark~\ref{r:KBtransOrder}).
The three-population network
\[
 \begin{aligned}
 \dot x_1&=-x_1+\Phi_a(x_1+\alpha x_2),\\
 \dot x_2&=-x_2+\Phi_a(x_2+\beta x_1x_3),\\
 \dot x_3&=-x_3+\Phi_a(x_3+\gamma x_2)
 \end{aligned}
\]
with positive couplings has overlapping minimal siphons
$\Sigma_{12}=\{1,2\}$ and $\Sigma_{23}=\{2,3\}$.
Their faces carry $(0,0,s_a)$ and $(s_a,0,0)$, where $s_a>0$
is the continued stable scalar equilibrium. Their invasion numbers are
\[
 R_{\Sigma_{12}}=a_1(1+\sqrt{\alpha\beta s_a}),\qquad
 R_{\Sigma_{23}}=a_1(1+\sqrt{\gamma\beta s_a}).
\]
They differ when $\alpha\ne\gamma$ and depend on the higher odds
coefficients through $s_a$.
These siphons share species $2$, whose production term $\beta x_1x_3$ enters both transverse matrices, while
both residents are independent of $\beta$. This \textbf{double boundary transcritical relay} takes the two residents
through transcritical bifurcations in turn by varying $\beta$ alone (Remark~\ref{r:KBrelay}, Figure~\ref{fig:relay}), which cannot
happen with disjoint siphons.
For sufficiently small positive couplings, Theorem~\ref{t:KBoverlapCusp}
constructs eight hyperbolic stable equilibria coexisting with an
interior cusp. Its positive cubic coefficient makes the middle
branch of the local unfolding stable, yielding at least nine stable
equilibria in a nearby open parameter region (Corollary~\ref{c:Nin}).
The cusp equilibrium itself is not counted among the eight.
A composite function, the scalar quartic energy on the resident coordinate plus a Perron-weighted
linear term on the absent coordinates, certifies an explicit subset of the basin of attraction of each boundary resident
(Proposition~\ref{p:ComBas}).
Earlier weak-coupling constructions \cite{Izhikevich1998,Boud26}
provide precedents for multiple stable states; the additional
constraints here preserve overlapping siphons and distinct explicit
boundary invasion thresholds.

\ideahead{5}
If the absent variables satisfy
\[
 \dot x_\Sigma\le(C_\Sigma-V_\Sigma)x_\Sigma,\qquad
 C_\Sigma\ge0,\qquad \rho(C_\Sigma V_\Sigma^{-1})<1,
\]
with $V_\Sigma$ a nonsingular $M$-matrix, then a constant weight
$w\gg0$ with $(C_\Sigma-V_\Sigma)^\top w\ll0$ certifies exponential
extinction (Theorem~\ref{t:Cmp}).
For a cooperative $C^2$ system on a compact invariant box with a
Hurwitz equilibrium $e$, Theorem~\ref{t:EqWgt} characterizes global
convergence to $e$ by a $C^1$ weight $v\gg0$ and $c>0$ satisfying
\[
 (\partial_xv)f+J^\top v\le-cv.
\]
The construction uses the variational flow
$\Psi(t,x)=\partial_x\varphi_t(x)$ of $f$, where
$\partial_t\Psi=J(\varphi_t(x))\Psi$ and $\Psi(0,x)=I$:
\[
 v_T(x)=\int_0^T\Psi(t,x)^\top q\,dt,\qquad q\gg0,
\]
for a sufficiently large finite $T$.
The contracting distance is
\[
 d_v(x,y)=\inf_{\gamma(0)=x,\ \gamma(1)=y}
 \int_0^1 v(\gamma(\tau))^\top|\gamma'(\tau)|\,d\tau,
\]
where the infimum is over piecewise $C^1$ paths in the box.
It need not equal $v(x)^\top|x-y|$.
Theorem~\ref{t:Wgt} gives the corresponding positive-weight converse
for Metzler variational equations over a compact flow.

The sufficient direction is established. For Metzler $J$, the inequality above is the
weighted $\ell^1$ case, $\Theta=\operatorname{Diag}(v)$, of Coogan's contraction condition
$\mu(\dot\Theta\Theta^{-1}+\Theta J\Theta^{-1})\le-c$ \cite{Coogan2019}; his separable Lyapunov
theorems and sum-of-squares search use weights $\theta_i(x_i)$.
The converse is also established for monotone systems: on a convex, forward and backward
invariant set, with bounded Jacobian, Kawano, Besselink and Cao
\cite[Theorem~4.3, Corollary~4.8]{KawanoBesselinkCao2020} prove that incremental exponential
stability is equivalent to a bounded, uniformly positive $C^1$ weight $v(x)$ satisfying this
inequality, and their proof uses the finite-horizon weight above
\cite[Appendix~B, (61)]{KawanoBesselinkCao2020}.
Theorem~\ref{t:Wgt} formulates this principle over a compact base flow: the base need not be
cooperative, $A$ need not be its Jacobian, and only forward evolution is used.
Theorem~\ref{t:EqWgt} is a compact-domain consequence, starting from global convergence to a
Hurwitz equilibrium on a forward invariant set. Existence of a polynomial weight on a
rectangular box is an approximation corollary. Separability here is in the tangent variables,
with $v_i$ depending on the full state.
Related converse results are in
\cite{LohmillerSlotine1998,ForniSepulchre2014,Giesl2015}.
The weights here also cover extinction. These results describe global convergence;
they are not prerequisites for the cusp constructions. Multistability
excludes contraction on a box containing the distinct equilibria.

Lyapunov--Volterra functions fit the same picture.
The function $\sum_ic_i(x_i-x_i^*-x_i^*\log(x_i/x_i^*))$ is a weighted sum of Fenchel--Bregman
divergences of the convex generator $x\log x-x$ of the exponential response
(Subsection~\ref{s:FenBre}). Its zero-reference limit gives the linear terms of the extinction
certificate of Theorem~\ref{t:Cmp}. The inverse-response primitive also enters the
response-integral energy of symmetric threshold networks (Proposition~\ref{p:SymEne}) and of the
scalar quartic odds equation. These divergences are nonnegative; the network
energy additionally contains a quadratic coupling term. Its decrease, and its local minimum at a
stable equilibrium, follow from the separate calculations given below.
Proposition~\ref{p:ComLya} combines a resident storage with such a linear term, and Proposition~\ref{p:LVboundary} applies this with a Lyapunov--Volterra storage to the three-population network.
\end{enumerate}

% \paragraph{Odds-polynomial responses and the minimum cusp degree.}
% Autoregulatory genetic models already combine saturating production,
% positive feedback and bistability; Isaacs et al.~\cite{Isaacs2003}
% provide a theoretical and experimental example.
% Multiple cusp bifurcations and multiple attracting states in weakly
% connected neural networks were studied by Izhikevich
% \cite{Izhikevich1998}.
% The specific contribution of
% Proposition~\ref{p:Qua}  is the exact minimum degree within the stated
% odds-polynomial class: degree four is necessary and sufficient for a
% stable origin and a positive nondegenerate cusp to coexist.
% The response zero set fixes the siphons, its onset slope determines
% the transverse production derivatives at a fixed resident state,
% and higher coefficients control the interior equilibrium geometry.
% Theorem~\ref{t:Net} then continues the cusp under coupling that
% preserves the prescribed boundary faces.
%
% \paragraph{Overlapping siphons and coexistence of stable states.}
% Persistence of hyperbolic equilibria under small perturbations is
% standard, and weakly coupled multistable networks already occur in
% \cite{Izhikevich1998}. Theorem~\ref{t:KBoverlapCusp} imposes additional
% boundary constraints on such a construction.
% The minimal siphons $\{1,2\}$ and $\{2,3\}$ remain fixed, their proper
% faces carry nonzero resident equilibria with distinct invasion numbers,
% and eight hyperbolic stable equilibria coexist with an interior cusp.
% Corollary~\ref{c:Nin} obtains a ninth stable equilibrium from its
% unfolding. This is an explicit realization of the general
% siphon-preservation rule, rather than a claim that weak coupling
% or multistability itself is new.
%
% \paragraph{Positive weights and converse contraction.}
% Coogan \cite{Coogan2019} develops contraction criteria using
% state-dependent weighted norms, while Kawano, Besselink and Cao
% \cite{KawanoBesselinkCao2020} prove converse results for monotone
% systems and positive linear time-varying systems.
% Theorem~\ref{t:EqWgt} belongs to this established line of work.
% For a cooperative system on a compact invariant box, it constructs
% positive contraction weights from global convergence to a Hurwitz
% equilibrium. Its role here is to turn convergence, including
% extinction, into an explicit variational-flow certificate.
% The siphon framework identifies the relevant boundary subsystems;
% the converse contraction principle has an independent literature..
% Keep the abstract's existing \ideaitem / \ideahead definitions unchanged.
% The six headings below are supplied by \ideahead, not copied manually.
\subsection{Main results}\label{s:mainRes}


The following enumeration follows the abstract; the text gives the main
formulas, their scope and relevant earlier results.
\begin{enumerate}
\ideahead{1}
The siphon identity and positivity give
\[
 J(0,y)=\begin{pmatrix}
 J_\Sigma^\perp(y)&0\\
 J_{\mathrm{mix},\Sigma}(y)&J_{\mathrm{tan},\Sigma}(y)
 \end{pmatrix},\qquad
 \begin{aligned}
 J_\Sigma^\perp&=\partial_zf_\Sigma(0,y),\\
 J_{\mathrm{mix},\Sigma}&=\partial_zf_{\Sigma^c}(0,y),\\
 J_{\mathrm{tan},\Sigma}&=\partial_yf_{\Sigma^c}(0,y).
 \end{aligned}
\]
The transverse block is Metzler even when the resident system is not
cooperative (Proposition~\ref{p:Tri}, \cite{AAHK,AA26}). For a regular
splitting $J_\Sigma^\perp=F_\Sigma-V_\Sigma$, with
$F_\Sigma\ge0$ and $V_\Sigma$ a nonsingular $M$-matrix,
\[
 s(J_\Sigma^\perp)<0
 \quad\Longleftrightarrow\quad
 R_\Sigma:=\rho(F_\Sigma V_\Sigma^{-1})<1
\]
is the next-generation matrix criterion \cite{Diek},\newline \cite{Van}.

For example, in the higher-order SIS model \eqref{eq:SIS}, $V_\Sigma=\operatorname{Diag}(\mu_i:i\in\Sigma)$.
For $i,j\in\Sigma$,
\[
 (F_\Sigma)_{ij}(y)=a_{ij}+
 \sum_{\substack{|E|\ge2\\E\cap\Sigma=\{j\}}}
 \beta_{iE}\prod_{k\in E\setminus\{j\}}y_k,
 \qquad J_\Sigma^\perp=F_\Sigma-V_\Sigma.
\]
Thus any term containing exactly one absent variable contributes to
invasion, regardless of its total degree. Terms vanishing together
with their first $z$-derivatives on the face leave both the resident
equations and the invasion matrix unchanged.
The scalar simplicial model \eqref{eq:simp} already exhibits bistability although its derivative at zero,
$\lambda-1$, is independent of $\lambda_\Delta$.

Theorem~\ref{t:siphMulti} constructs any prescribed finite number
of stable positive equilibria while preserving every siphon and every
transverse Jacobian $-V_\Sigma$. It uses nonnegative polynomials
$P_i$ and a smooth response vanishing on $(-\infty,0]$ and
positive on $(0,\infty)$. Its support condition is $V^{-1}c\gg0$,
where $c_i=b_i$ if $P_i\not\equiv0$ and $c_i=0$ otherwise.
The response depends on the prescribed count; no common slope bound
is imposed. The response is flat at onset, so all production derivatives on
siphon faces vanish: the transverse Jacobians are specifically
$-V_\Sigma$, rather than arbitrary prescribed matrices. This controls all boundary faces simultaneously, extending
the separation illustrated by the scalar contagion model. Siphon-based
persistence and boundary analysis have precedents in \cite{AdLS}.
The face property alone does not imply interior cooperativity:
$\dot z_1=-z_1z_2$, $\dot z_2=-z_2$ has a Metzler Jacobian at zero
but $\partial_{z_2}f_1=-z_1<0$ for $z_1>0$.

\ideahead{4}
For $\dot x=-Vx+\Phi(\beta Wx)$, assume a synaptic matrix $W\ge0$, $\beta\ge0$ and
that $\Phi\ge0$ is smooth on the attained arguments with $0\le \Phi'\le L$,
$L>0$. Theorem~\ref{t:KBsharp} compares
\[
 \bar R_0=\rho(\beta WV^{-1})
\]
with the reciprocal $1/L$ of the slope bound.
It extends the next-generation number: $J(x)=F(x)-V$ with
$F(x)=\beta\operatorname{Diag}(\Phi'(\beta Wx))W\le\beta LW$, so the classical number
$R(x_*)=\rho(F(x_*)V^{-1})$, which decides the stability of one equilibrium $x_*$ \cite{Van}, satisfies
$R(x_*)\le L\bar R_0$ for all equilibria and all responses in the class (Remark~\ref{r:NGMg}).
The transition is
\[
 \begin{cases}
 \bar R_0<\dfrac1L: & \parbox[t]{.7\linewidth}{every response in the class gives a unique globally exponentially attracting equilibrium, since $J(x)\le-V+\beta LW$;}\\[12pt]
 \bar R_0>\dfrac1L: & \parbox[t]{.7\linewidth}{for each $m\ge2$, some response in the class gives at least $m$ locally exponentially stable equilibria, including a choice with $\Phi>0$ and $\Phi'>0$ on $\mathbb R$ and all equilibria strictly positive.}
 \end{cases}
\]
No assertion at equality is needed. The response may depend on $m$;
this is not a count for one fixed response or polynomial degree.

Positive-feedback multistability is studied generally in
\cite{AngeliSontag2004}, \cite{AngeliFerrellSontag2004}.
The contraction below the bound is known: it is the case $B=I$
of the multivariable Lur'e theorem of \cite{Davydov}, for general $V$,
and \cite{Boud26} has the case $V=I$, $W\mathbf1=\mathbf1$; see also
\cite{Coogan2019,KawanoBesselinkCao2020}. The contribution here is its
constructive sharpness: with $V$, $W$, $L$ fixed and one common scalar response, every
prescribed finite stable count is attained above the bound, without irreducibility.

The construction principle, plateaus of the response on which $J=-V$, has a predecessor.
For $\dot u_i=-d_iu_i+\sum_jw_{ij}f_j(u_j)+I_i$ with piecewise linear
activations with $\ell+1$ plateaus, \cite[Theorem~3]{WangLuChen2010} give exactly
$(2\ell+1)^n$ equilibria, of which $(\ell+1)^n$ are locally stable, and estimate their
attraction basins; multilevel sigmoid activations are in their Section~5. Their condition (35)
requires self-feedback domination, which for a common activation implies
$w_{ii}>\ell\sum_{j\ne i}|w_{ij}|$. Thus it
cannot supply arbitrary counts on a fixed network with cross-coupling, whereas
Theorem~\ref{t:KBsharp} needs only $\bar R_0>1/L$, and asserts a lower bound rather than an exact count.
For example, take $V=I$, $W=\begin{pmatrix}0&1\\1&0\end{pmatrix}$ and $\beta L>1$. With
$y=\beta Wx$ the system becomes $\dot y=-y+\beta W\Phi(y)$, which has zero diagonal weights, so
(35) cannot hold, while Theorem~\ref{t:KBsharp} gives every finite count.
Theorem~\ref{t:siphMulti} uses the same mechanism; its addition is that $V$ and the
polynomials $P_i$ stay fixed, the siphon lattice and every transverse Jacobian are preserved,
and $V^{-1}c\gg0$ is the exact feasibility condition.

Betteti, Baggio, Bullo and Zampieri \cite{Betteti2025} construct symmetric firing-rate networks
whose prescribed two-level memory patterns are stable equilibria. Their formulas prescribe
response values at two argument levels, and their stability test controls the derivatives at those
levels, a precedent for the value-and-slope construction used here. Their synaptic matrix is
designed from the patterns, whereas Theorem~\ref{t:KBsharp} fixes the removal and synaptic matrices, dimension
and response slope bound and realizes every finite stable count by varying one common response.
Their construction does not impose the simultaneous preservation of siphons and transverse
matrices established in Theorem~\ref{t:siphMulti}. In their symmetric setting, an energy function
gives convergence to the set of equilibria. Its integral term is the convex generator of the
Fenchel--Bregman divergence of Subsection~\ref{s:FenBre}, which also recovers Volterra functions.

For an explicit three-population network, Theorem~\ref{t:KBsharp} is compatible with the
$2$-contraction criterion of \cite{OfirOvseevichMargaliot2024}. As a consequence of the two
results, arbitrarily many stable equilibria can coexist while every trajectory converges to an
equilibrium.

Equilibria of \eqref{eq:KB} are the fixed points of $T=V^{-1}\Pi$, and $DT(x_*)=V^{-1}\partial_x\Pi(x_*)$ is similar to the
next-generation matrix $\partial_x\Pi(x_*)V^{-1}$. For linear $P_i$, $T$ is a two-layer nonnegative neural network in the sense of
\cite{piotrowski2024fixed}. Their results apply when $T$ is monotone and weakly scalable, as for concave responses: the fixed point
is then unique and positive if $T$ is scalable, and the positive fixed points form an interval of a ray under a primitivity
condition. The responses of Theorems~\ref{t:KBsharp} and~\ref{t:siphMulti} and Proposition~\ref{p:Qua} are therefore not concave
where they are evaluated. Their certificate $T^m(0)\gg0$ excludes every nonempty siphon, since $T$ maps each siphon face into itself.

Linear-threshold networks have diagonal removal and clipped affine production, which is only locally Lipschitz.
\cite{Retnaraj2026} study a timescale family, with production $\tau^{-1}[Dx+\tau((W-D)x+h)]_0^1$, whose equilibrium set is independent
of $\tau$. Its invariant siphon faces are independent of $\tau$ as well: on $x_i=0$, the production vanishes exactly when
$((W-D)x+h)_i\le0$. Smooth bifurcation criteria require additional regularity at the clipping thresholds.

Unlike Theorem~\ref{t:siphMulti}, Theorem~\ref{t:KBsharp} fixes a slope bound
and linear $P_i$; its full-support construction has no nonempty
siphon. Neither theorem requires the other, and neither constructs
a cusp. The explicit cusp results use separate calculations and
continuation arguments.

\ifdefined\withSpatial
\ideahead{spat}
Add diffusion to \eqref{eq:KB}, $x_t=\mathcal D\Delta_\zeta x+\Pi(x)-Vx$ with $\zeta\in\mathbb R^n$, and linearize at a resident equilibrium on the face of a siphon
$\Sigma$. The absent variables obey $z_t=\mathcal D_\Sigma\Delta_\zeta z+J_\Sigma^\perp z$, with $J_\Sigma^\perp=F_\Sigma-V_\Sigma$, and the ansatz
$z=e^{-\lambda(e\cdot\zeta-ct)}q$, with $e$ a unit vector, gives the speed of exponential modes normal to $e$. In principle this speed depends on $e$; since the Laplacian assumes diffusion independent of the direction, it does not, and it is the speed of Proposition~\ref{t:spread}. The threshold is again next-generation: $c_*>0$ exactly when
$R_\Sigma>1$, and $c_*\sim2\sqrt{\langle\mathcal D\rangle s(J_\Sigma^\perp)}$ near it (Remark~\ref{r:movNGM}). The speed has the form of the
Freidlin--G\"artner formula. It concerns the linearized equation only: the flat-onset responses of Theorem~\ref{t:siphMulti} have
$R_\Sigma=0$ and no such speed, and with bistability a front can have positive speed although $R_\Sigma<1$.
For threshold responses, Theorem~\ref{t:SpdBd} carries Theorem~\ref{t:KBsharp} over to diffusion: below $1/L$ every solution converges to the
one equilibrium and there are no fronts, above it the $m$ stable states persist for every $\mathcal D$, and no response in $\mathcal C_L$ with $\Phi(0)=0$ spreads faster
than $\bar c(L)$, which we conjecture degree-one odds responses attain (Remark~\ref{r:SpdAtt}). For odds responses only $a_1$ enters the speed at a fixed resident
(Corollary~\ref{c:OddSpd}). The relay switches on two linear speeds, but numerically a front invades $E_{12}$ already while it is stable (Remark~\ref{r:RelSpd}), and the shared species $x_2$, absent from both residents, ignites where they meet (Remark~\ref{r:OvlIgn}, Figure~\ref{fig:overFront}). Every homogeneous Hurwitz equilibrium of a cooperative system stays stable under arbitrary diagonal diffusion, so all the constructed counts of stable equilibria survive, including the nine of Corollary~\ref{c:Nin} (Corollary~\ref{c:DifSta}); reductions with group-compatible diffusion, extinction certificates, global convergence on a box and, on bounded domains, homogeneous cusps extend as well (Corollary~\ref{c:DifAll}). Integrated Lyapunov functions acquire a diffusion term, which has a sign for Lyapunov--Volterra functions but not in general for network energies, and the composite basin certificate of Proposition~\ref{p:ComBas} extends after a change of storage (Proposition~\ref{p:DifLya}). Section~\ref{s:spaSta} recalls the Evans function and treats the stability of these fronts; above $c_*$ an exponential weight stabilizes the state ahead (Proposition~\ref{p:WeiSpe}).
\fi

\ideahead{3}
Partition the species into nonempty groups $C_1,\ldots,C_{n_G}$ and
let $H_{ia}=1$ for $i\in C_a$ and zero otherwise. Then $x=Hy$
assigns the common value $y_a$ to group $C_a$.
Assume $f(Hy;\theta)=H\widehat f(y;\theta)$ for all admissible
$y,\theta$. The group-constant states are invariant, and
\[
 J(Hy;\theta)H=H\widehat J(y;\theta),\qquad
 s(J(Hy;\theta))=s(\widehat J(y;\theta))
 \quad\text{if }J(Hy;\theta)\text{ is Metzler}.
\]
Theorem~\ref{t:Spe} transfers local exponential equilibrium stability
without irreducibility. A nondegenerate saddle-node or cusp also lifts
when the full Jacobian has a simple zero eigenvalue and all remaining
eigenvalues have negative real parts. These reductions support the
homogeneous and two-population calculations in
Theorems~\ref{t:KBhom} and~\ref{t:KBexchange}; the homogeneous cusp of a unimodal density
gives the hysteresis loop of Figure~\ref{fig:hyst}. For a symmetric response,
$\Phi(-q)=1-\Phi(q)$, Theorem~\ref{t:cuspAsym} gives the exact
asymmetry interval $|\eta|<\eta_*(\beta)$ on which both two-population cusps remain
admissible and nondegenerate.
Invariant synchrony subspaces and balanced partitions are established
in \cite{SGP03,golubitsky2005patterns}; here the cooperative spectral
equality supplies the stability test. Exact quotient reduction for row-stochastic
threshold networks, $WH=HB$, is in \cite{Boud26}, where a lifted equilibrium $Hy^*$ is
stable under the sufficient condition $\beta\max_aq_a<1$, with
$q_a=\Phi'(h+\beta(By^*)_a)$; the spectral equality replaces it by the exact test
$\beta\rho(\operatorname{Diag}(q)B)<1$. Invariance preserves initial
equality within groups, but does not imply synchronization from
arbitrary initial states.

\ideahead{odds}
For $\varphi(q)=\sum_{k=1}^m a_kq^k$, with $a_k\ge0$ and $a_m>0$,
put $\Phi(q)=\varphi(q)/(1+\varphi(q))$ for $q\ge0$ and $\Phi(q)=0$ for $q\le0$.
The zero set is fixed, and $\Phi'(0+)=a_1$.
For fixed $V$, $b_i>0$ and nonnegative polynomials $P_i$,
Proposition~\ref{p:Odd} preserves the siphon lattice and gives
\[
 J_\Sigma^\perp(y)=-V_\Sigma
 +a_1\operatorname{Diag}(b_i:i\in\Sigma)D_zP_\Sigma(0,y).
\]
Higher coefficients enter an invasion threshold only indirectly,
through the resident equilibrium. For $\dot x=-x+\phi(x)/(1+\phi(x))$,
with $\phi(0)=0$ and nonnegative polynomial coefficients, Proposition~\ref{p:Qua} combines
the exact minimum degree, four, for a positive nondegenerate cusp with a hyperbolically
stable origin, with an attaining construction.
Its reference polynomial is
\[
 \phi(q)=\frac{342q+368q^2+128q^3+1024q^4}{343}.
\]
For $\dot x=-x+\Phi(x)$, the equilibria $0$ and $1/2$ are hyperbolically
stable, while $1/8$ is a nondegenerate cusp with $f'''(1/8)=36/49$,
generically unfolded by $(a_1,a_2)$. More generally, the root parameters
$0<r<s$, $3r+s\le1$, give
\[
 c=[(1-r)^3(1-s)]^{-1},\qquad
 1-a_1=cr^3s,\qquad f'''(r)=6cr(s-r)(1-r).
\]
Within this family, the margin $1-a_1$ has supremum $1/81$, approached only as the cusp
degenerates. This bounds the uncoupled origin's stability margin,
not the coupling range of a continued cusp.
Theorem~\ref{t:Net} continues the cusp and the two stable equilibria
under coupling that preserves singleton siphons and a transversely
stable boundary equilibrium.

Bistability and hysteresis under rational genetic feedback are established:
\cite{Hasty2001} derive
\begin{equation}\label{eq:Hasty}
 \dot x=m\frac{1+x^2+\alpha\sigma_1x^4}{1+x^2+\sigma_1x^4+\sigma_1\sigma_2x^6}-\gamma_xx
\end{equation}
by eliminating fast dimerization and promoter-binding reactions, with hysteresis as the degradation
rate varies; see also \cite{Isaacs2003}. Its production at zero is $m>0$, so the origin is not an equilibrium, and its numerator
and denominator are different polynomials; its $x^4$ term therefore has a different role from
the degree of $\phi$ here.
Earlier biological feedback models already exhibit cusps associated
with rational response functions. In particular, van Hoek and
Hogeweg \cite[Eq.~(21)]{vanHoek} locate a cusp by imposing
a triple root on the polynomial equilibrium equation.
Izhikevich \cite[Theorem~1]{Izhikevich1998} derives coupled cubic
normal forms for weakly coupled systems whose uncoupled components
already have cusps. These results precede our use of polynomial
equilibrium equations and coupling near a cusp.

We have not found the minimum-degree result in the preceding
literature; it concerns the displayed odds family, not rational genetic responses in general. The local cusp criterion is standard \cite{kuznetsov2004};
the additional result is its explicit realization subject to these
positivity and boundary constraints.

\ideahead{overl}
Responses positive at every finite argument, such as Gaussian and logistic distribution functions,
exclude nonempty siphons (Remark~\ref{r:KBnosiph}). For zero-onset responses,
Proposition~\ref{p:InvF} gives the criterion:
for each siphon $\Sigma$, every production monomial for a species
$i\in\Sigma$ must contain at least one variable from $\Sigma$,
and $V_{\Sigma\Sigma^c}=0$. For a singleton, this forces its own
variable to divide each production monomial (Remark~\ref{r:Mul}).
For linear $P_i$, minimal siphons are the source strongly connected components of the
graph of $W$ and $V$ (Remark~\ref{r:KBsiphgraph}); for higher-order $P_i$ they can overlap
(Example~\ref{e:KBoverlap}). The degree of $P_i$ in the absent variables and the onset order
of $\Phi_i$ fix the order of production along transverse rays (Remark~\ref{r:KBtransOrder}).
The three-population network
\[
 \begin{aligned}
 \dot x_1&=-x_1+\Phi_a(x_1+\alpha x_2),\\
 \dot x_2&=-x_2+\Phi_a(x_2+\beta x_1x_3),\\
 \dot x_3&=-x_3+\Phi_a(x_3+\gamma x_2)
 \end{aligned}
\]
with positive couplings has overlapping minimal siphons
$\Sigma_{12}=\{1,2\}$ and $\Sigma_{23}=\{2,3\}$.
Their faces carry $(0,0,s_a)$ and $(s_a,0,0)$, where $s_a>0$
is the continued stable scalar equilibrium. Their invasion numbers are
\[
 R_{\Sigma_{12}}=a_1(1+\sqrt{\alpha\beta s_a}),\qquad
 R_{\Sigma_{23}}=a_1(1+\sqrt{\gamma\beta s_a}).
\]
They differ when $\alpha\ne\gamma$ and depend on the higher odds
coefficients through $s_a$.
These siphons share species $2$, whose production term $\beta x_1x_3$ enters both transverse matrices, while
both residents are independent of $\beta$. This \textbf{double boundary transcritical relay} takes the two residents
through transcritical bifurcations in turn by varying $\beta$ alone (Remark~\ref{r:KBrelay}, Figure~\ref{fig:relay}), which cannot
happen with disjoint siphons.
For sufficiently small positive couplings, Theorem~\ref{t:KBoverlapCusp}
constructs eight hyperbolic stable equilibria coexisting with an
interior cusp. Its positive cubic coefficient makes the middle
branch of the local unfolding stable, yielding at least nine stable
equilibria in a nearby open parameter region (Corollary~\ref{c:Nin}).
The cusp equilibrium itself is not counted among the eight.
A composite function, the scalar quartic energy on the resident coordinate plus a Perron-weighted
linear term on the absent coordinates, certifies an explicit subset of the basin of attraction of each boundary resident
(Proposition~\ref{p:ComBas}).
Earlier weak-coupling constructions \cite{Izhikevich1998,Boud26}
provide precedents for multiple stable states; the additional
constraints here preserve overlapping siphons and distinct explicit
boundary invasion thresholds.

\ideahead{5}
If the absent variables satisfy
\[
 \dot x_\Sigma\le(C_\Sigma-V_\Sigma)x_\Sigma,\qquad
 C_\Sigma\ge0,\qquad \rho(C_\Sigma V_\Sigma^{-1})<1,
\]
with $V_\Sigma$ a nonsingular $M$-matrix, then a constant weight
$w\gg0$ with $(C_\Sigma-V_\Sigma)^\top w\ll0$ certifies exponential
extinction (Theorem~\ref{t:Cmp}).
For a cooperative $C^2$ system on a compact invariant box with a
Hurwitz equilibrium $e$, Theorem~\ref{t:EqWgt} characterizes global
convergence to $e$ by a $C^1$ weight $v\gg0$ and $c>0$ satisfying
\[
 (\partial_xv)f+J^\top v\le-cv.
\]
The construction uses the variational flow
$\Psi(t,x)=\partial_x\varphi_t(x)$ of $f$, where
$\partial_t\Psi=J(\varphi_t(x))\Psi$ and $\Psi(0,x)=I$:
\[
 v_T(x)=\int_0^T\Psi(t,x)^\top q\,dt,\qquad q\gg0,
\]
for a sufficiently large finite $T$.
The contracting distance is
\[
 d_v(x,y)=\inf_{\gamma(0)=x,\ \gamma(1)=y}
 \int_0^1 v(\gamma(\tau))^\top|\gamma'(\tau)|\,d\tau,
\]
where the infimum is over piecewise $C^1$ paths in the box.
It need not equal $v(x)^\top|x-y|$.
Theorem~\ref{t:Wgt} gives the corresponding positive-weight converse
for Metzler variational equations over a compact flow.

The sufficient direction is established. For Metzler $J$, the inequality above is the
weighted $\ell^1$ case, $\Theta=\operatorname{Diag}(v)$, of Coogan's contraction condition
$\mu(\dot\Theta\Theta^{-1}+\Theta J\Theta^{-1})\le-c$ \cite{Coogan2019}; his separable Lyapunov
theorems and sum-of-squares search use weights $\theta_i(x_i)$.
The converse is also established for monotone systems: on a convex, forward and backward
invariant set, with bounded Jacobian, Kawano, Besselink and Cao
\cite[Theorem~4.3, Corollary~4.8]{KawanoBesselinkCao2020} prove that incremental exponential
stability is equivalent to a bounded, uniformly positive $C^1$ weight $v(x)$ satisfying this
inequality, and their proof uses the finite-horizon weight above
\cite[Appendix~B, (61)]{KawanoBesselinkCao2020}.
Theorem~\ref{t:Wgt} formulates this principle over a compact base flow: the base need not be
cooperative, $A$ need not be its Jacobian, and only forward evolution is used.
Theorem~\ref{t:EqWgt} is a compact-domain consequence, starting from global convergence to a
Hurwitz equilibrium on a forward invariant set. Existence of a polynomial weight on a
rectangular box is an approximation corollary. Separability here is in the tangent variables,
with $v_i$ depending on the full state.
Related converse results are in
\cite{LohmillerSlotine1998,ForniSepulchre2014,Giesl2015}.
The weights here also cover extinction. These results describe global convergence;
they are not prerequisites for the cusp constructions. Multistability
excludes contraction on a box containing the distinct equilibria.

Lyapunov--Volterra functions fit the same picture.
The function $\sum_ic_i(x_i-x_i^*-x_i^*\log(x_i/x_i^*))$ is a weighted sum of Fenchel--Bregman
divergences of the convex generator $x\log x-x$ of the exponential response
(Subsection~\ref{s:FenBre}). Its zero-reference limit gives the linear terms of the extinction
certificate of Theorem~\ref{t:Cmp}. The inverse-response primitive also enters the
response-integral energy of symmetric threshold networks (Proposition~\ref{p:SymEne}) and of the
scalar quartic odds equation. These divergences are nonnegative; the network
energy additionally contains a quadratic coupling term. Its decrease, and its local minimum at a
stable equilibrium, follow from the separate calculations given below.
Proposition~\ref{p:ComLya} combines a resident storage with such a linear term, and Proposition~\ref{p:LVboundary} applies this with a Lyapunov--Volterra storage to the three-population network.
\end{enumerate}

% \paragraph{Odds-polynomial responses and the minimum cusp degree.}
% Autoregulatory genetic models already combine saturating production,
% positive feedback and bistability; Isaacs et al.~\cite{Isaacs2003}
% provide a theoretical and experimental example.
% Multiple cusp bifurcations and multiple attracting states in weakly
% connected neural networks were studied by Izhikevich
% \cite{Izhikevich1998}.
% The specific contribution of
% Proposition~\ref{p:Qua}  is the exact minimum degree within the stated
% odds-polynomial class: degree four is necessary and sufficient for a
% stable origin and a positive nondegenerate cusp to coexist.
% The response zero set fixes the siphons, its onset slope determines
% the transverse production derivatives at a fixed resident state,
% and higher coefficients control the interior equilibrium geometry.
% Theorem~\ref{t:Net} then continues the cusp under coupling that
% preserves the prescribed boundary faces.
%
% \paragraph{Overlapping siphons and coexistence of stable states.}
% Persistence of hyperbolic equilibria under small perturbations is
% standard, and weakly coupled multistable networks already occur in
% \cite{Izhikevich1998}. Theorem~\ref{t:KBoverlapCusp} imposes additional
% boundary constraints on such a construction.
% The minimal siphons $\{1,2\}$ and $\{2,3\}$ remain fixed, their proper
% faces carry nonzero resident equilibria with distinct invasion numbers,
% and eight hyperbolic stable equilibria coexist with an interior cusp.
% Corollary~\ref{c:Nin} obtains a ninth stable equilibrium from its
% unfolding. This is an explicit realization of the general
% siphon-preservation rule, rather than a claim that weak coupling
% or multistability itself is new.
%
% \paragraph{Positive weights and converse contraction.}
% Coogan \cite{Coogan2019} develops contraction criteria using
% state-dependent weighted norms, while Kawano, Besselink and Cao
% \cite{KawanoBesselinkCao2020} prove converse results for monotone
% systems and positive linear time-varying systems.
% Theorem~\ref{t:EqWgt} belongs to this established line of work.
% For a cooperative system on a compact invariant box, it constructs
% positive contraction weights from global convergence to a Hurwitz
% equilibrium. Its role here is to turn convergence, including
% extinction, into an explicit variational-flow certificate.
% The siphon framework identifies the relevant boundary subsystems;
% the converse contraction principle has an independent literature..
% Keep the abstract's existing \ideaitem / \ideahead definitions unchanged.
% The six headings below are supplied by \ideahead, not copied manually.
\subsection{Main results}\label{s:mainRes}


The following enumeration follows the abstract; the text gives the main
formulas, their scope and relevant earlier results.
\begin{enumerate}
\ideahead{1}
The siphon identity and positivity give
\[
 J(0,y)=\begin{pmatrix}
 J_\Sigma^\perp(y)&0\\
 J_{\mathrm{mix},\Sigma}(y)&J_{\mathrm{tan},\Sigma}(y)
 \end{pmatrix},\qquad
 \begin{aligned}
 J_\Sigma^\perp&=\partial_zf_\Sigma(0,y),\\
 J_{\mathrm{mix},\Sigma}&=\partial_zf_{\Sigma^c}(0,y),\\
 J_{\mathrm{tan},\Sigma}&=\partial_yf_{\Sigma^c}(0,y).
 \end{aligned}
\]
The transverse block is Metzler even when the resident system is not
cooperative (Proposition~\ref{p:Tri}, \cite{AAHK,AA26}). For a regular
splitting $J_\Sigma^\perp=F_\Sigma-V_\Sigma$, with
$F_\Sigma\ge0$ and $V_\Sigma$ a nonsingular $M$-matrix,
\[
 s(J_\Sigma^\perp)<0
 \quad\Longleftrightarrow\quad
 R_\Sigma:=\rho(F_\Sigma V_\Sigma^{-1})<1
\]
is the next-generation matrix criterion \cite{Diek},\newline \cite{Van}.

For example, in the higher-order SIS model \eqref{eq:SIS}, $V_\Sigma=\operatorname{Diag}(\mu_i:i\in\Sigma)$.
For $i,j\in\Sigma$,
\[
 (F_\Sigma)_{ij}(y)=a_{ij}+
 \sum_{\substack{|E|\ge2\\E\cap\Sigma=\{j\}}}
 \beta_{iE}\prod_{k\in E\setminus\{j\}}y_k,
 \qquad J_\Sigma^\perp=F_\Sigma-V_\Sigma.
\]
Thus any term containing exactly one absent variable contributes to
invasion, regardless of its total degree. Terms vanishing together
with their first $z$-derivatives on the face leave both the resident
equations and the invasion matrix unchanged.
The scalar simplicial model \eqref{eq:simp} already exhibits bistability although its derivative at zero,
$\lambda-1$, is independent of $\lambda_\Delta$.

Theorem~\ref{t:siphMulti} constructs any prescribed finite number
of stable positive equilibria while preserving every siphon and every
transverse Jacobian $-V_\Sigma$. It uses nonnegative polynomials
$P_i$ and a smooth response vanishing on $(-\infty,0]$ and
positive on $(0,\infty)$. Its support condition is $V^{-1}c\gg0$,
where $c_i=b_i$ if $P_i\not\equiv0$ and $c_i=0$ otherwise.
The response depends on the prescribed count; no common slope bound
is imposed. The response is flat at onset, so all production derivatives on
siphon faces vanish: the transverse Jacobians are specifically
$-V_\Sigma$, rather than arbitrary prescribed matrices. This controls all boundary faces simultaneously, extending
the separation illustrated by the scalar contagion model. Siphon-based
persistence and boundary analysis have precedents in \cite{AdLS}.
The face property alone does not imply interior cooperativity:
$\dot z_1=-z_1z_2$, $\dot z_2=-z_2$ has a Metzler Jacobian at zero
but $\partial_{z_2}f_1=-z_1<0$ for $z_1>0$.

\ideahead{4}
For $\dot x=-Vx+\Phi(\beta Wx)$, assume a synaptic matrix $W\ge0$, $\beta\ge0$ and
that $\Phi\ge0$ is smooth on the attained arguments with $0\le \Phi'\le L$,
$L>0$. Theorem~\ref{t:KBsharp} compares
\[
 \bar R_0=\rho(\beta WV^{-1})
\]
with the reciprocal $1/L$ of the slope bound.
It extends the next-generation number: $J(x)=F(x)-V$ with
$F(x)=\beta\operatorname{Diag}(\Phi'(\beta Wx))W\le\beta LW$, so the classical number
$R(x_*)=\rho(F(x_*)V^{-1})$, which decides the stability of one equilibrium $x_*$ \cite{Van}, satisfies
$R(x_*)\le L\bar R_0$ for all equilibria and all responses in the class (Remark~\ref{r:NGMg}).
The transition is
\[
 \begin{cases}
 \bar R_0<\dfrac1L: & \parbox[t]{.7\linewidth}{every response in the class gives a unique globally exponentially attracting equilibrium, since $J(x)\le-V+\beta LW$;}\\[12pt]
 \bar R_0>\dfrac1L: & \parbox[t]{.7\linewidth}{for each $m\ge2$, some response in the class gives at least $m$ locally exponentially stable equilibria, including a choice with $\Phi>0$ and $\Phi'>0$ on $\mathbb R$ and all equilibria strictly positive.}
 \end{cases}
\]
No assertion at equality is needed. The response may depend on $m$;
this is not a count for one fixed response or polynomial degree.

Positive-feedback multistability is studied generally in
\cite{AngeliSontag2004}, \cite{AngeliFerrellSontag2004}.
The contraction below the bound is known: it is the case $B=I$
of the multivariable Lur'e theorem of \cite{Davydov}, for general $V$,
and \cite{Boud26} has the case $V=I$, $W\mathbf1=\mathbf1$; see also
\cite{Coogan2019,KawanoBesselinkCao2020}. The contribution here is its
constructive sharpness: with $V$, $W$, $L$ fixed and one common scalar response, every
prescribed finite stable count is attained above the bound, without irreducibility.

The construction principle, plateaus of the response on which $J=-V$, has a predecessor.
For $\dot u_i=-d_iu_i+\sum_jw_{ij}f_j(u_j)+I_i$ with piecewise linear
activations with $\ell+1$ plateaus, \cite[Theorem~3]{WangLuChen2010} give exactly
$(2\ell+1)^n$ equilibria, of which $(\ell+1)^n$ are locally stable, and estimate their
attraction basins; multilevel sigmoid activations are in their Section~5. Their condition (35)
requires self-feedback domination, which for a common activation implies
$w_{ii}>\ell\sum_{j\ne i}|w_{ij}|$. Thus it
cannot supply arbitrary counts on a fixed network with cross-coupling, whereas
Theorem~\ref{t:KBsharp} needs only $\bar R_0>1/L$, and asserts a lower bound rather than an exact count.
For example, take $V=I$, $W=\begin{pmatrix}0&1\\1&0\end{pmatrix}$ and $\beta L>1$. With
$y=\beta Wx$ the system becomes $\dot y=-y+\beta W\Phi(y)$, which has zero diagonal weights, so
(35) cannot hold, while Theorem~\ref{t:KBsharp} gives every finite count.
Theorem~\ref{t:siphMulti} uses the same mechanism; its addition is that $V$ and the
polynomials $P_i$ stay fixed, the siphon lattice and every transverse Jacobian are preserved,
and $V^{-1}c\gg0$ is the exact feasibility condition.

Betteti, Baggio, Bullo and Zampieri \cite{Betteti2025} construct symmetric firing-rate networks
whose prescribed two-level memory patterns are stable equilibria. Their formulas prescribe
response values at two argument levels, and their stability test controls the derivatives at those
levels, a precedent for the value-and-slope construction used here. Their synaptic matrix is
designed from the patterns, whereas Theorem~\ref{t:KBsharp} fixes the removal and synaptic matrices, dimension
and response slope bound and realizes every finite stable count by varying one common response.
Their construction does not impose the simultaneous preservation of siphons and transverse
matrices established in Theorem~\ref{t:siphMulti}. In their symmetric setting, an energy function
gives convergence to the set of equilibria. Its integral term is the convex generator of the
Fenchel--Bregman divergence of Subsection~\ref{s:FenBre}, which also recovers Volterra functions.

For an explicit three-population network, Theorem~\ref{t:KBsharp} is compatible with the
$2$-contraction criterion of \cite{OfirOvseevichMargaliot2024}. As a consequence of the two
results, arbitrarily many stable equilibria can coexist while every trajectory converges to an
equilibrium.

Equilibria of \eqref{eq:KB} are the fixed points of $T=V^{-1}\Pi$, and $DT(x_*)=V^{-1}\partial_x\Pi(x_*)$ is similar to the
next-generation matrix $\partial_x\Pi(x_*)V^{-1}$. For linear $P_i$, $T$ is a two-layer nonnegative neural network in the sense of
\cite{piotrowski2024fixed}. Their results apply when $T$ is monotone and weakly scalable, as for concave responses: the fixed point
is then unique and positive if $T$ is scalable, and the positive fixed points form an interval of a ray under a primitivity
condition. The responses of Theorems~\ref{t:KBsharp} and~\ref{t:siphMulti} and Proposition~\ref{p:Qua} are therefore not concave
where they are evaluated. Their certificate $T^m(0)\gg0$ excludes every nonempty siphon, since $T$ maps each siphon face into itself.

Linear-threshold networks have diagonal removal and clipped affine production, which is only locally Lipschitz.
\cite{Retnaraj2026} study a timescale family, with production $\tau^{-1}[Dx+\tau((W-D)x+h)]_0^1$, whose equilibrium set is independent
of $\tau$. Its invariant siphon faces are independent of $\tau$ as well: on $x_i=0$, the production vanishes exactly when
$((W-D)x+h)_i\le0$. Smooth bifurcation criteria require additional regularity at the clipping thresholds.

Unlike Theorem~\ref{t:siphMulti}, Theorem~\ref{t:KBsharp} fixes a slope bound
and linear $P_i$; its full-support construction has no nonempty
siphon. Neither theorem requires the other, and neither constructs
a cusp. The explicit cusp results use separate calculations and
continuation arguments.

\ifdefined\withSpatial
\ideahead{spat}
Add diffusion to \eqref{eq:KB}, $x_t=\mathcal D\Delta_\zeta x+\Pi(x)-Vx$ with $\zeta\in\mathbb R^n$, and linearize at a resident equilibrium on the face of a siphon
$\Sigma$. The absent variables obey $z_t=\mathcal D_\Sigma\Delta_\zeta z+J_\Sigma^\perp z$, with $J_\Sigma^\perp=F_\Sigma-V_\Sigma$, and the ansatz
$z=e^{-\lambda(e\cdot\zeta-ct)}q$, with $e$ a unit vector, gives the speed of exponential modes normal to $e$. In principle this speed depends on $e$; since the Laplacian assumes diffusion independent of the direction, it does not, and it is the speed of Proposition~\ref{t:spread}. The threshold is again next-generation: $c_*>0$ exactly when
$R_\Sigma>1$, and $c_*\sim2\sqrt{\langle\mathcal D\rangle s(J_\Sigma^\perp)}$ near it (Remark~\ref{r:movNGM}). The speed has the form of the
Freidlin--G\"artner formula. It concerns the linearized equation only: the flat-onset responses of Theorem~\ref{t:siphMulti} have
$R_\Sigma=0$ and no such speed, and with bistability a front can have positive speed although $R_\Sigma<1$.
For threshold responses, Theorem~\ref{t:SpdBd} carries Theorem~\ref{t:KBsharp} over to diffusion: below $1/L$ every solution converges to the
one equilibrium and there are no fronts, above it the $m$ stable states persist for every $\mathcal D$, and no response in $\mathcal C_L$ with $\Phi(0)=0$ spreads faster
than $\bar c(L)$, which we conjecture degree-one odds responses attain (Remark~\ref{r:SpdAtt}). For odds responses only $a_1$ enters the speed at a fixed resident
(Corollary~\ref{c:OddSpd}). The relay switches on two linear speeds, but numerically a front invades $E_{12}$ already while it is stable (Remark~\ref{r:RelSpd}), and the shared species $x_2$, absent from both residents, ignites where they meet (Remark~\ref{r:OvlIgn}, Figure~\ref{fig:overFront}). Every homogeneous Hurwitz equilibrium of a cooperative system stays stable under arbitrary diagonal diffusion, so all the constructed counts of stable equilibria survive, including the nine of Corollary~\ref{c:Nin} (Corollary~\ref{c:DifSta}); reductions with group-compatible diffusion, extinction certificates, global convergence on a box and, on bounded domains, homogeneous cusps extend as well (Corollary~\ref{c:DifAll}). Integrated Lyapunov functions acquire a diffusion term, which has a sign for Lyapunov--Volterra functions but not in general for network energies, and the composite basin certificate of Proposition~\ref{p:ComBas} extends after a change of storage (Proposition~\ref{p:DifLya}). Section~\ref{s:spaSta} recalls the Evans function and treats the stability of these fronts; above $c_*$ an exponential weight stabilizes the state ahead (Proposition~\ref{p:WeiSpe}).
\fi

\ideahead{3}
Partition the species into nonempty groups $C_1,\ldots,C_{n_G}$ and
let $H_{ia}=1$ for $i\in C_a$ and zero otherwise. Then $x=Hy$
assigns the common value $y_a$ to group $C_a$.
Assume $f(Hy;\theta)=H\widehat f(y;\theta)$ for all admissible
$y,\theta$. The group-constant states are invariant, and
\[
 J(Hy;\theta)H=H\widehat J(y;\theta),\qquad
 s(J(Hy;\theta))=s(\widehat J(y;\theta))
 \quad\text{if }J(Hy;\theta)\text{ is Metzler}.
\]
Theorem~\ref{t:Spe} transfers local exponential equilibrium stability
without irreducibility. A nondegenerate saddle-node or cusp also lifts
when the full Jacobian has a simple zero eigenvalue and all remaining
eigenvalues have negative real parts. These reductions support the
homogeneous and two-population calculations in
Theorems~\ref{t:KBhom} and~\ref{t:KBexchange}; the homogeneous cusp of a unimodal density
gives the hysteresis loop of Figure~\ref{fig:hyst}. For a symmetric response,
$\Phi(-q)=1-\Phi(q)$, Theorem~\ref{t:cuspAsym} gives the exact
asymmetry interval $|\eta|<\eta_*(\beta)$ on which both two-population cusps remain
admissible and nondegenerate.
Invariant synchrony subspaces and balanced partitions are established
in \cite{SGP03,golubitsky2005patterns}; here the cooperative spectral
equality supplies the stability test. Exact quotient reduction for row-stochastic
threshold networks, $WH=HB$, is in \cite{Boud26}, where a lifted equilibrium $Hy^*$ is
stable under the sufficient condition $\beta\max_aq_a<1$, with
$q_a=\Phi'(h+\beta(By^*)_a)$; the spectral equality replaces it by the exact test
$\beta\rho(\operatorname{Diag}(q)B)<1$. Invariance preserves initial
equality within groups, but does not imply synchronization from
arbitrary initial states.

\ideahead{odds}
For $\varphi(q)=\sum_{k=1}^m a_kq^k$, with $a_k\ge0$ and $a_m>0$,
put $\Phi(q)=\varphi(q)/(1+\varphi(q))$ for $q\ge0$ and $\Phi(q)=0$ for $q\le0$.
The zero set is fixed, and $\Phi'(0+)=a_1$.
For fixed $V$, $b_i>0$ and nonnegative polynomials $P_i$,
Proposition~\ref{p:Odd} preserves the siphon lattice and gives
\[
 J_\Sigma^\perp(y)=-V_\Sigma
 +a_1\operatorname{Diag}(b_i:i\in\Sigma)D_zP_\Sigma(0,y).
\]
Higher coefficients enter an invasion threshold only indirectly,
through the resident equilibrium. For $\dot x=-x+\phi(x)/(1+\phi(x))$,
with $\phi(0)=0$ and nonnegative polynomial coefficients, Proposition~\ref{p:Qua} combines
the exact minimum degree, four, for a positive nondegenerate cusp with a hyperbolically
stable origin, with an attaining construction.
Its reference polynomial is
\[
 \phi(q)=\frac{342q+368q^2+128q^3+1024q^4}{343}.
\]
For $\dot x=-x+\Phi(x)$, the equilibria $0$ and $1/2$ are hyperbolically
stable, while $1/8$ is a nondegenerate cusp with $f'''(1/8)=36/49$,
generically unfolded by $(a_1,a_2)$. More generally, the root parameters
$0<r<s$, $3r+s\le1$, give
\[
 c=[(1-r)^3(1-s)]^{-1},\qquad
 1-a_1=cr^3s,\qquad f'''(r)=6cr(s-r)(1-r).
\]
Within this family, the margin $1-a_1$ has supremum $1/81$, approached only as the cusp
degenerates. This bounds the uncoupled origin's stability margin,
not the coupling range of a continued cusp.
Theorem~\ref{t:Net} continues the cusp and the two stable equilibria
under coupling that preserves singleton siphons and a transversely
stable boundary equilibrium.

Bistability and hysteresis under rational genetic feedback are established:
\cite{Hasty2001} derive
\begin{equation}\label{eq:Hasty}
 \dot x=m\frac{1+x^2+\alpha\sigma_1x^4}{1+x^2+\sigma_1x^4+\sigma_1\sigma_2x^6}-\gamma_xx
\end{equation}
by eliminating fast dimerization and promoter-binding reactions, with hysteresis as the degradation
rate varies; see also \cite{Isaacs2003}. Its production at zero is $m>0$, so the origin is not an equilibrium, and its numerator
and denominator are different polynomials; its $x^4$ term therefore has a different role from
the degree of $\phi$ here.
Earlier biological feedback models already exhibit cusps associated
with rational response functions. In particular, van Hoek and
Hogeweg \cite[Eq.~(21)]{vanHoek} locate a cusp by imposing
a triple root on the polynomial equilibrium equation.
Izhikevich \cite[Theorem~1]{Izhikevich1998} derives coupled cubic
normal forms for weakly coupled systems whose uncoupled components
already have cusps. These results precede our use of polynomial
equilibrium equations and coupling near a cusp.

We have not found the minimum-degree result in the preceding
literature; it concerns the displayed odds family, not rational genetic responses in general. The local cusp criterion is standard \cite{kuznetsov2004};
the additional result is its explicit realization subject to these
positivity and boundary constraints.

\ideahead{overl}
Responses positive at every finite argument, such as Gaussian and logistic distribution functions,
exclude nonempty siphons (Remark~\ref{r:KBnosiph}). For zero-onset responses,
Proposition~\ref{p:InvF} gives the criterion:
for each siphon $\Sigma$, every production monomial for a species
$i\in\Sigma$ must contain at least one variable from $\Sigma$,
and $V_{\Sigma\Sigma^c}=0$. For a singleton, this forces its own
variable to divide each production monomial (Remark~\ref{r:Mul}).
For linear $P_i$, minimal siphons are the source strongly connected components of the
graph of $W$ and $V$ (Remark~\ref{r:KBsiphgraph}); for higher-order $P_i$ they can overlap
(Example~\ref{e:KBoverlap}). The degree of $P_i$ in the absent variables and the onset order
of $\Phi_i$ fix the order of production along transverse rays (Remark~\ref{r:KBtransOrder}).
The three-population network
\[
 \begin{aligned}
 \dot x_1&=-x_1+\Phi_a(x_1+\alpha x_2),\\
 \dot x_2&=-x_2+\Phi_a(x_2+\beta x_1x_3),\\
 \dot x_3&=-x_3+\Phi_a(x_3+\gamma x_2)
 \end{aligned}
\]
with positive couplings has overlapping minimal siphons
$\Sigma_{12}=\{1,2\}$ and $\Sigma_{23}=\{2,3\}$.
Their faces carry $(0,0,s_a)$ and $(s_a,0,0)$, where $s_a>0$
is the continued stable scalar equilibrium. Their invasion numbers are
\[
 R_{\Sigma_{12}}=a_1(1+\sqrt{\alpha\beta s_a}),\qquad
 R_{\Sigma_{23}}=a_1(1+\sqrt{\gamma\beta s_a}).
\]
They differ when $\alpha\ne\gamma$ and depend on the higher odds
coefficients through $s_a$.
These siphons share species $2$, whose production term $\beta x_1x_3$ enters both transverse matrices, while
both residents are independent of $\beta$. This \textbf{double boundary transcritical relay} takes the two residents
through transcritical bifurcations in turn by varying $\beta$ alone (Remark~\ref{r:KBrelay}, Figure~\ref{fig:relay}), which cannot
happen with disjoint siphons.
For sufficiently small positive couplings, Theorem~\ref{t:KBoverlapCusp}
constructs eight hyperbolic stable equilibria coexisting with an
interior cusp. Its positive cubic coefficient makes the middle
branch of the local unfolding stable, yielding at least nine stable
equilibria in a nearby open parameter region (Corollary~\ref{c:Nin}).
The cusp equilibrium itself is not counted among the eight.
A composite function, the scalar quartic energy on the resident coordinate plus a Perron-weighted
linear term on the absent coordinates, certifies an explicit subset of the basin of attraction of each boundary resident
(Proposition~\ref{p:ComBas}).
Earlier weak-coupling constructions \cite{Izhikevich1998,Boud26}
provide precedents for multiple stable states; the additional
constraints here preserve overlapping siphons and distinct explicit
boundary invasion thresholds.

\ideahead{5}
If the absent variables satisfy
\[
 \dot x_\Sigma\le(C_\Sigma-V_\Sigma)x_\Sigma,\qquad
 C_\Sigma\ge0,\qquad \rho(C_\Sigma V_\Sigma^{-1})<1,
\]
with $V_\Sigma$ a nonsingular $M$-matrix, then a constant weight
$w\gg0$ with $(C_\Sigma-V_\Sigma)^\top w\ll0$ certifies exponential
extinction (Theorem~\ref{t:Cmp}).
For a cooperative $C^2$ system on a compact invariant box with a
Hurwitz equilibrium $e$, Theorem~\ref{t:EqWgt} characterizes global
convergence to $e$ by a $C^1$ weight $v\gg0$ and $c>0$ satisfying
\[
 (\partial_xv)f+J^\top v\le-cv.
\]
The construction uses the variational flow
$\Psi(t,x)=\partial_x\varphi_t(x)$ of $f$, where
$\partial_t\Psi=J(\varphi_t(x))\Psi$ and $\Psi(0,x)=I$:
\[
 v_T(x)=\int_0^T\Psi(t,x)^\top q\,dt,\qquad q\gg0,
\]
for a sufficiently large finite $T$.
The contracting distance is
\[
 d_v(x,y)=\inf_{\gamma(0)=x,\ \gamma(1)=y}
 \int_0^1 v(\gamma(\tau))^\top|\gamma'(\tau)|\,d\tau,
\]
where the infimum is over piecewise $C^1$ paths in the box.
It need not equal $v(x)^\top|x-y|$.
Theorem~\ref{t:Wgt} gives the corresponding positive-weight converse
for Metzler variational equations over a compact flow.

The sufficient direction is established. For Metzler $J$, the inequality above is the
weighted $\ell^1$ case, $\Theta=\operatorname{Diag}(v)$, of Coogan's contraction condition
$\mu(\dot\Theta\Theta^{-1}+\Theta J\Theta^{-1})\le-c$ \cite{Coogan2019}; his separable Lyapunov
theorems and sum-of-squares search use weights $\theta_i(x_i)$.
The converse is also established for monotone systems: on a convex, forward and backward
invariant set, with bounded Jacobian, Kawano, Besselink and Cao
\cite[Theorem~4.3, Corollary~4.8]{KawanoBesselinkCao2020} prove that incremental exponential
stability is equivalent to a bounded, uniformly positive $C^1$ weight $v(x)$ satisfying this
inequality, and their proof uses the finite-horizon weight above
\cite[Appendix~B, (61)]{KawanoBesselinkCao2020}.
Theorem~\ref{t:Wgt} formulates this principle over a compact base flow: the base need not be
cooperative, $A$ need not be its Jacobian, and only forward evolution is used.
Theorem~\ref{t:EqWgt} is a compact-domain consequence, starting from global convergence to a
Hurwitz equilibrium on a forward invariant set. Existence of a polynomial weight on a
rectangular box is an approximation corollary. Separability here is in the tangent variables,
with $v_i$ depending on the full state.
Related converse results are in
\cite{LohmillerSlotine1998,ForniSepulchre2014,Giesl2015}.
The weights here also cover extinction. These results describe global convergence;
they are not prerequisites for the cusp constructions. Multistability
excludes contraction on a box containing the distinct equilibria.

Lyapunov--Volterra functions fit the same picture.
The function $\sum_ic_i(x_i-x_i^*-x_i^*\log(x_i/x_i^*))$ is a weighted sum of Fenchel--Bregman
divergences of the convex generator $x\log x-x$ of the exponential response
(Subsection~\ref{s:FenBre}). Its zero-reference limit gives the linear terms of the extinction
certificate of Theorem~\ref{t:Cmp}. The inverse-response primitive also enters the
response-integral energy of symmetric threshold networks (Proposition~\ref{p:SymEne}) and of the
scalar quartic odds equation. These divergences are nonnegative; the network
energy additionally contains a quadratic coupling term. Its decrease, and its local minimum at a
stable equilibrium, follow from the separate calculations given below.
Proposition~\ref{p:ComLya} combines a resident storage with such a linear term, and Proposition~\ref{p:LVboundary} applies this with a Lyapunov--Volterra storage to the three-population network.
\end{enumerate}

% \paragraph{Odds-polynomial responses and the minimum cusp degree.}
% Autoregulatory genetic models already combine saturating production,
% positive feedback and bistability; Isaacs et al.~\cite{Isaacs2003}
% provide a theoretical and experimental example.
% Multiple cusp bifurcations and multiple attracting states in weakly
% connected neural networks were studied by Izhikevich
% \cite{Izhikevich1998}.
% The specific contribution of
% Proposition~\ref{p:Qua}  is the exact minimum degree within the stated
% odds-polynomial class: degree four is necessary and sufficient for a
% stable origin and a positive nondegenerate cusp to coexist.
% The response zero set fixes the siphons, its onset slope determines
% the transverse production derivatives at a fixed resident state,
% and higher coefficients control the interior equilibrium geometry.
% Theorem~\ref{t:Net} then continues the cusp under coupling that
% preserves the prescribed boundary faces.
%
% \paragraph{Overlapping siphons and coexistence of stable states.}
% Persistence of hyperbolic equilibria under small perturbations is
% standard, and weakly coupled multistable networks already occur in
% \cite{Izhikevich1998}. Theorem~\ref{t:KBoverlapCusp} imposes additional
% boundary constraints on such a construction.
% The minimal siphons $\{1,2\}$ and $\{2,3\}$ remain fixed, their proper
% faces carry nonzero resident equilibria with distinct invasion numbers,
% and eight hyperbolic stable equilibria coexist with an interior cusp.
% Corollary~\ref{c:Nin} obtains a ninth stable equilibrium from its
% unfolding. This is an explicit realization of the general
% siphon-preservation rule, rather than a claim that weak coupling
% or multistability itself is new.
%
% \paragraph{Positive weights and converse contraction.}
% Coogan \cite{Coogan2019} develops contraction criteria using
% state-dependent weighted norms, while Kawano, Besselink and Cao
% \cite{KawanoBesselinkCao2020} prove converse results for monotone
% systems and positive linear time-varying systems.
% Theorem~\ref{t:EqWgt} belongs to this established line of work.
% For a cooperative system on a compact invariant box, it constructs
% positive contraction weights from global convergence to a Hurwitz
% equilibrium. Its role here is to turn convergence, including
% extinction, into an explicit variational-flow certificate.
% The siphon framework identifies the relevant boundary subsystems;
% the converse contraction principle has an independent literature.